\documentclass[12pt,a4paper]{article}

\usepackage{fontspec}
\setmainfont{Times New Roman}
\usepackage{amsmath,amssymb,bm}
\usepackage[top=2.5cm,bottom=2.5cm,left=3cm,right=2.5cm]{geometry}
\usepackage{setspace}
\onehalfspacing
\usepackage{indentfirst}
\setlength{\parindent}{1.5em}
\setlength{\parskip}{0.35em}
\setlength{\emergencystretch}{1em}
\usepackage{graphicx}
\usepackage{booktabs}
\usepackage{multirow}
\usepackage{array}
\usepackage{tabularx}
\usepackage{longtable}
\usepackage{xcolor}
\usepackage{float}
\usepackage{caption}
\usepackage{subcaption}
\usepackage{tikz}
\usetikzlibrary{arrows.meta,positioning,fit,calc,shapes.geometric}
\usepackage{adjustbox}
\usepackage{pdflscape}
\usepackage{hyperref}
\usepackage{url}
\usepackage{cite}
\newcolumntype{P}[1]{>{\raggedright\arraybackslash}p{#1}}
\newcolumntype{Y}{>{\raggedright\arraybackslash}X}

\hypersetup{
  colorlinks=true,
  linkcolor=blue!55!black,
  citecolor=blue!55!black,
  urlcolor=blue!55!black,
  pdftitle={Donor-Slot Set Transformer cho nhận dạng người đóng góp và phát hiện ngoài-panel trong hỗn hợp DNA STR},
  pdfauthor={Nguyễn Vũ Mạnh, Lê Sỹ Nguyên, Nguyễn Đắc Quang, Trịnh Hoàng Giang, Nguyễn Quốc Dũng}
}

\renewcommand\thesection{\Roman{section}}
\renewcommand\thesubsection{\arabic{section}.\arabic{subsection}}
\renewcommand\thesubsubsection{\arabic{section}.\arabic{subsection}.\arabic{subsubsection}}
\renewcommand{\contentsname}{Mục lục}
\renewcommand{\figurename}{Hình}
\renewcommand{\tablename}{Bảng}
\renewcommand{\refname}{Tài liệu tham khảo}

\definecolor{peakblue}{RGB}{43,108,176}
\definecolor{locusgreen}{RGB}{47,133,90}
\definecolor{profileorange}{RGB}{221,107,32}
\definecolor{headpurple}{RGB}{128,90,213}
\definecolor{lightgray}{RGB}{245,247,250}

\tikzset{
  block/.style={draw,rounded corners=2pt,align=center,minimum height=1.15cm,
    text width=2.35cm,fill=blue!6,font=\small},
  smallblock/.style={draw,rounded corners=2pt,align=center,minimum height=0.95cm,
    text width=1.85cm,fill=gray!7,font=\scriptsize},
  line/.style={-{Latex[length=2.4mm]},thick},
  groupbox/.style={draw,dashed,rounded corners=4pt,inner sep=7pt}
}

\begin{document}

\begin{titlepage}
  \centering
  \vspace*{1.0cm}

  {\large\textbf{TRƯỜNG ĐẠI HỌC BÁCH KHOA HÀ NỘI}}\\[0.25cm]
  {\normalsize Trường Công nghệ Thông tin và Truyền thông}\\[1.6cm]

  \rule{\linewidth}{1.4pt}\\[0.5cm]
  {\LARGE\bfseries Donor-Slot Set Transformer: nhận dạng người đóng góp}\\[0.25cm]
  {\LARGE\bfseries và phát hiện ngoài-panel trong hỗn hợp DNA STR}\\[0.5cm]
  \rule{\linewidth}{1.4pt}\\[1.4cm]

  {\large\textbf{Báo cáo nghiên cứu}}\\[1.3cm]

  \begin{tabular}{>{\bfseries}rl}
    Thành viên nhóm: & Nguyễn Vũ Mạnh \\[0.25cm]
                     & Lê Sỹ Nguyên \\[0.25cm]
                     & Nguyễn Đắc Quang \\[0.25cm]
                     & Trịnh Hoàng Giang \\[0.25cm]
                     & Nguyễn Quốc Dũng
  \end{tabular}

  \vfill
  {\large Hà Nội, tháng 6 năm 2026}
\end{titlepage}

\tableofcontents
\newpage

\section*{Tóm tắt}
\addcontentsline{toc}{section}{Tóm tắt}

Phân tích hỗn hợp DNA là một trong những nhiệm vụ khó nhất của giám định di truyền pháp y vì tín hiệu điện di quan sát được có thể đồng thời chứa allele của nhiều người, kèm stutter, nhiễu nền, allele dropout và chênh lệch lớn về tỷ lệ đóng góp. Hai câu hỏi trung tâm khi diễn giải là những người nào trong một tập tham chiếu đã góp DNA, và liệu có người đóng góp nào nằm ngoài tập tham chiếu hay không. Nghiên cứu này trình bày một mô hình duy nhất, gọi là Donor-Slot Set Transformer (DS-ST), với một bộ mã hóa tập dùng chung phục vụ hai nhiệm vụ chính: phân loại người đóng góp trong panel (CLS) và từ chối open-set (reject). Số người đóng góp (NOC) không phải mục tiêu chính mà là một đại lượng suy ra từ hồ sơ xác suất CLS bằng một bộ đếm hậu nghiệm; trong quá trình phát triển, đại lượng suy ra này được phát hiện đạt độ chính xác rất cao nên cũng được báo cáo như một kết quả phụ. Mỗi peak được biểu diễn bằng một token tám trường gồm chỉ số locus, allele, log chiều cao và các đặc trưng quan hệ; bảy trường số liên tục được mã hóa bằng periodic numerical embedding còn locus bằng embedding học được, sau đó chuỗi token đi qua hai khối ISAB++ với 32 inducing point, bốn attention head, SetNorm, clean-path residual và attention sigmoid không cạnh tranh tại bước inducing point. Từ biểu diễn dùng chung, một bộ giải mã khe có điều kiện kiểu gen — khởi tạo theo CoSA, tinh chỉnh theo GSANet, định tuyến bằng Sinkhorn optimal transport theo MESH và gate tồn tại theo AdaSlot — trả về xác suất từng donor cho nhiệm vụ CLS, và một đầu từ chối trên nhánh pooling không lọc đảm nhiệm reject; số người được suy ra từ hồ sơ xác suất bằng một bộ đếm hậu nghiệm. Hai đầu phụ gồm attribution theo peak và hồi quy tỷ lệ hỗn hợp định hình biểu diễn khi huấn luyện, kèm một bước tái xếp hạng theo tỷ lệ hỗn hợp khi suy luận để làm sắc nét thứ hạng donor. Dữ liệu được xây dựng từ PROVEDIt, cơ sở dữ liệu công khai hơn 27.000 hỗn hợp DNA pháp y \cite{provedit2018,proveditweb}: nhánh đếm dùng các hồ sơ GlobalFiler thực, còn nhánh nhận dạng huấn luyện trên hỗn hợp bán-tổng hợp trộn từ hồ sơ nguồn đơn thực và được đánh giá trên hỗn hợp thực với tổ hợp người đóng góp tách biệt hoàn toàn khỏi tập huấn luyện. Ở nhiệm vụ nhận dạng, DS-ST đạt exact match \(0,962\), Macro F1 \(0,987\) và reject AUROC \(0,999\), dẫn đầu các baseline cổ điển; ở bài toán đếm suy ra, DS-ST đạt Macro F1 \(0,9444\pm0,0042\) trên GroupKFold năm phần, cao hơn XGBoost \(0,9230\pm0,0209\), CNN một chiều \(0,8854\pm0,0100\) và Forensic Transformer \(0,8846\pm0,0131\). Ablation thành phần xác nhận giám sát đặc quyền nâng trần xếp hạng còn bộ giải mã khe nâng chỉ số triển khai, trong khi nhiều thành phần được khảo sát rồi loại bỏ vì không cải thiện.

\medskip
\noindent\textbf{Từ khóa:} hỗn hợp DNA; số lượng người đóng góp; STR; Set Transformer; Deep Sets; học sâu pháp y; giải trộn DNA; PROVEDIt.

\section*{Các đóng góp chính}
\addcontentsline{toc}{section}{Các đóng góp chính}

Đóng góp thứ nhất là một mô hình duy nhất, Donor-Slot Set Transformer, với một bộ mã hóa tập dùng chung phục vụ \emph{ba} nhiệm vụ pháp y: đếm số người đóng góp (NOC), phân loại người đóng góp trong panel (CLS) và từ chối open-set (reject). Thiết kế này tránh xây dựng các pipeline tách rời cho từng nhiệm vụ, đồng thời bảo toàn tính bất biến đối với hoán vị peak vốn là yêu cầu tự nhiên của dữ liệu dạng tập hợp \cite{zaheer2017,lee2019}.

Đóng góp thứ hai là bộ mã hóa tập của DS-ST với token tám trường, periodic numerical embedding, hai khối ISAB++, attention sigmoid không cạnh tranh, SetNorm, clean-path residual, 32 inducing point và lọc khả thi cùng tách tập riêng/chia sẻ. Chính bộ mã hóa này, khi ghép với một đầu đếm năm lớp, giải quyết nhiệm vụ NOC; khi ghép với bộ giải mã khe, giải quyết nhiệm vụ CLS.

Đóng góp thứ ba là bộ giải mã khe có điều kiện kiểu gen cho nhiệm vụ CLS trên 45 donor tham chiếu, cùng một đầu từ chối tách rời cho nhiệm vụ reject. Các khái niệm Conditional Slot Attention, Guided Slot Attention, MESH và Adaptive Slot Attention được điều chỉnh sang bối cảnh pháp y, trong đó mỗi khe đại diện cho một donor, ma trận Sinkhorn biểu diễn phân bổ bằng chứng peak đến donor và gate Gumbel-Sigmoid biểu diễn khả năng donor tồn tại trong hỗn hợp \cite{ma2024cosa,lee2024gsa,zhang2023mesh,fan2024adaslot}.

Đóng góp thứ tư là quy trình đánh giá chống rò rỉ: nhánh đếm dùng GroupKFold theo danh tính mixture trên dữ liệu thực, còn nhánh nhận dạng huấn luyện trên hỗn hợp bán-tổng hợp và được đánh giá trên hỗn hợp thực với tổ hợp người đóng góp tách biệt hoàn toàn. Nghiên cứu báo cáo benchmark với XGBoost, MLP, hồi quy logistic, CNN, kNN, TabPFN và Forensic Transformer; phân tích ablation thành phần cho DS-ST; đồng thời xây dựng dữ liệu cho GlobalFiler, Identifiler Plus và PowerPlex 16 để chuẩn bị đánh giá khác miền.

\section{Giới thiệu}

\subsection{Bối cảnh giám định DNA hỗn hợp}

DNA profiling dựa trên short tandem repeat đã trở thành một trụ cột của giám định sinh học kể từ khi các mẫu lặp đa hình được sử dụng cho nhận dạng cá thể vào thập niên 1980 \cite{jeffreys1985}. Trong quy trình capillary electrophoresis hiện đại, các vùng STR được khuếch đại, phân tách theo kích thước và biểu diễn thành các peak huỳnh quang. Một hồ sơ nguồn đơn thường có không quá hai allele thật tại mỗi locus autosomal, nhưng hồ sơ hỗn hợp có thể chứa nhiều hơn hai allele, các peak có độ cao chênh lệch mạnh và các artefact phát sinh từ phản ứng PCR hoặc thiết bị. PROVEDIt được xây dựng để hỗ trợ kiểm chứng các phương pháp phân tích trong bối cảnh này và hiện được Rutgers mô tả là cơ sở dữ liệu mở gồm hơn 27.000 hỗn hợp DNA có ý nghĩa pháp y \cite{provedit2018,proveditweb}.

Trong một vụ việc thực tế, mẫu DNA có thể được thu từ máu, tinh dịch, tế bào biểu mô hoặc dấu vết tiếp xúc. Số người để lại vật chất sinh học thường không được biết trước. Việc ước lượng Number of Contributors vì thế là bước đầu của nhiều quy trình diễn giải. NOCIt, STRmix, EuroForMix và các hệ thống probabilistic genotyping khác đều cần mô hình hóa một hoặc nhiều giả thuyết về số contributor; giả thuyết sai có thể làm thay đổi likelihood ratio và cách đánh giá sức nặng của chứng cứ \cite{haned2011,bleka2016,bright2016,coble2019}. NIST nhấn mạnh rằng khả năng phát hiện DNA từ lượng tế bào ngày càng nhỏ làm tăng số trường hợp xuất hiện hỗn hợp phức tạp và khiến việc phân biệt nguồn trở nên khó hơn \cite{nist2024}.

Phương pháp đơn giản nhất thường dựa trên maximum allele count. Nếu locus có \(A_{\max}\) allele, một cận dưới tự nhiên cho số người là \(\lceil A_{\max}/2\rceil\). Quy tắc này dễ giải thích nhưng mất thông tin về chiều cao peak, tần suất allele, dropout và stutter. Haned và cộng sự đã xây dựng cách tiếp cận xác suất so sánh giả thuyết hai người, ba người và các số contributor khác bằng xác suất allele quan sát được \cite{haned2011}. Alfonse và cộng sự đánh giá NOCIt, maximum likelihood và maximum allele count trên 100 mixture phức tạp, cho thấy việc suy luận trở nên khó hơn rõ rệt ở chế độ low template \cite{alfonse2017}. Những kết quả này cho thấy NOC không chỉ là bài toán đếm số allele mà là bài toán suy luận từ tín hiệu nhiễu, thiếu và chồng lấn.

\subsection{Hai nhiệm vụ chính và bài toán đếm suy ra}

Nhiệm vụ chính thứ nhất là phân loại người đóng góp (CLS): xác định tập contributor từ một panel tham chiếu có \(C\) donor, với nhãn đa nhãn \(\bm{y}^{\mathrm{id}}\in\{0,1\}^{C}\) mà tổng số phần tử bằng một chính là số người thật khi mọi contributor đều nằm trong panel. Đây là nhiệm vụ khó nhất vì mô hình phải phân biệt các donor chia sẻ allele và phải tìm bằng chứng của contributor thiểu số khi peak của họ bị che bởi contributor chính. Nghiên cứu sequencing của Phan và cộng sự dùng 27 STR và 94 SNP từ sáu cá thể để dự đoán cá thể trong hỗn hợp, minh họa tiềm năng của học sâu cho nhận dạng contributor nhưng cũng cho thấy thiết lập panel nhỏ và dữ liệu MPS khác đáng kể so với điện di mao quản \cite{phan2021}.

Nhiệm vụ chính thứ hai là từ chối open-set (reject): vì panel tham chiếu là tập đóng, một hỗn hợp thực có thể chứa người đóng góp không thuộc panel, và mô hình cần phát hiện trường hợp này thay vì gán nhầm vào các donor đã biết. Đây là một bài toán phân loại nhị phân riêng (có/không có contributor ngoài panel), được giải bằng một đầu và một nhánh gộp tách rời, và là điều kiện cần để triển khai an toàn trong thực tế.

Số người đóng góp (NOC) không được xem là một nhiệm vụ chính mà là một đại lượng \emph{suy ra} từ chất lượng của CLS: khi hồ sơ xác suất donor đã sắc nét, số người có thể đọc ra trực tiếp từ hồ sơ đó. Cụ thể, bộ đếm triển khai là một rừng ngẫu nhiên hậu nghiệm chạy trên hồ sơ xác suất CLS, được hỗ trợ bởi một số hạng giám sát đếm nhẹ khi huấn luyện. Trong quá trình phát triển, chúng tôi quan sát thấy đại lượng suy ra này đạt độ chính xác đếm rất cao (xem phần thử nghiệm), dù không phải mục tiêu chính; vì vậy nghiên cứu báo cáo nó như một kết quả phụ đáng chú ý và là hướng còn nhiều dư địa cải tiến. Ba đầu ra hỗ trợ nhau — NOC cung cấp số \(k\) để giải mã top \(k\) cho CLS, còn reject bảo vệ CLS khỏi giả định panel đóng — nhưng vẫn dùng đầu ra và hàm mất mát riêng trên một bộ mã hóa chung.

\subsection{Động lực của biểu diễn dạng tập peak}

Véc tơ phẳng theo cột allele có ưu điểm đơn giản và phù hợp với XGBoost hoặc TabPFN, nhưng cấu trúc của nó phụ thuộc bộ kit và không biểu diễn trực tiếp quan hệ giữa các peak trong cùng locus. CNN một chiều áp dụng convolution dọc theo véc tơ cột, song tính lân cận giữa hai cột không nhất thiết tương ứng với quan hệ sinh học. Forensic Transformer cải thiện điều này bằng cách tạo một token cho mỗi locus, nhưng việc nén toàn bộ peak của một locus thành một véc tơ trước attention có thể làm mất bằng chứng yếu của contributor thiểu số. Dữ liệu STR vì thế phù hợp hơn với một biểu diễn dạng \emph{tập peak}: mỗi peak là một phần tử độc lập và chỉ số locus được đưa vào chính token thay vì dùng để gộp peak, nhờ đó quan hệ giữa mọi peak — trong cùng locus lẫn giữa các locus — được học bằng attention mà không peak yếu nào bị trung bình hóa sớm.

Deep Sets chứng minh rằng một lớp hàm bất biến hoán vị có thể được biểu diễn dưới dạng \(\rho(\sum_i\phi(\bm{x}_i))\) trong các điều kiện thích hợp \cite{zaheer2017}. Set Transformer mở rộng ý tưởng này bằng attention, cho phép mô hình hóa tương tác giữa các phần tử trong tập và dùng inducing points để giảm độ phức tạp từ bậc hai xuống xấp xỉ \(O(NM)\), với \(N\) phần tử và \(M\) inducing points \cite{lee2019}. Kiến trúc thực tế của đồ án đi theo hướng thứ hai: mỗi peak là một phần tử của tập, chỉ số locus được đưa vào embedding, còn quan hệ giữa mọi peak được học qua ISAB++. Thiết kế không dựng tensor cố định theo từng locus và không dùng một Cross Locus Transformer riêng; cấu trúc locus được bảo toàn trong chính token và trong bảng tra genotype.

\subsection{Khoảng trống nghiên cứu}

PACE sử dụng 1.405 hồ sơ không mô phỏng chứa từ một đến bốn contributor, 15 locus Identifiler và đạt khoảng 97\% accuracy trên tập kiểm thử, nhưng phụ thuộc đặc trưng được lựa chọn và không mô hình hóa trực tiếp tập peak \cite{marciano2017,alamoudi2022}. RFC19 sử dụng 19 đặc trưng random forest trên 590 hồ sơ PowerPlex Fusion 6C và đạt accuracy được báo cáo khoảng 83,3\% \cite{benschop2019,kruijver2021}. TAWSEEM dùng MLP và báo cáo 97\% accuracy trên kịch bản lớn nhất gồm 6.000 profile từ bốn multiplex, nhưng đầu vào có tới 6.000 đặc trưng phẳng và cách chia dữ liệu khác với GroupKFold theo mixture identity \cite{alamoudi2022}. deepNoC mô phỏng 100.000 profile và đạt 89\% accuracy cho một đến mười contributor, sau đó fine tune bằng vài trăm profile phòng thí nghiệm \cite{taylor2024deepnoc}. Các công trình này cung cấp nền tảng quan trọng nhưng vẫn để ngỏ câu hỏi về một backbone có cấu trúc peak và locus rõ ràng, có thể tái sử dụng đồng thời cho NOC và contributor identification.

Một khoảng trống khác là tính minh bạch và khả năng tái lập. Veldhuis và cộng sự phát triển giải thích phản thực tế cho dự đoán NOC, trong đó explanation phải tạo ra hồ sơ thay đổi đủ thực tế thay vì chỉ thay đổi tùy ý một đặc trưng \cite{veldhuis2022}. Barash và cộng sự tổng quan các ứng dụng machine learning trong forensic DNA và nhấn mạnh nhu cầu kiểm chứng, khả năng diễn giải và hiểu biết liên ngành \cite{barash2024}. DNANet tiếp tục hướng mở bằng U Net cho allele calling; mô hình đạt F1 bằng 0,971 trên case data chưa thấy khi huấn luyện, 0,982 trên research data theo annotation chuyên gia và 0,962 trên allele donor thật trong phiên bản xuất bản \cite{dewit2025}. Mặc dù DNANet không giải quyết trực tiếp NOC, kết quả của nó cho thấy kiến trúc phổ biến, dữ liệu thực và mã nguồn mở có thể hạ thấp rào cản triển khai AI pháp y.

\section{Các nghiên cứu liên quan}

Lĩnh vực ước lượng số người đóng góp DNA hỗn hợp đã trải qua ba thế hệ phương pháp kế tiếp nhau: từ quy tắc đếm allele dựa trên luật, qua machine learning với đặc trưng thủ công, đến deep learning trên dữ liệu quy mô lớn hoặc mô phỏng. Bảng~\ref{tab:related_summary} tổng hợp ba công trình tiêu biểu đại diện cho ba thế hệ đó; các phần tiếp theo trình bày chi tiết nội dung, phương pháp và đóng góp của từng nhóm công trình, cũng như các nền tảng học sâu mà DS-ST kế thừa.

\begin{table}[H]
\centering
\caption{Tổng hợp ba nghiên cứu tiêu biểu về ước lượng Number of Contributors trong hỗn hợp DNA STR.}
\label{tab:related_summary}
\adjustbox{max width=\linewidth}{%
\small
\begin{tabular}{@{}P{3.2cm}lP{3.6cm}P{5.2cm}P{5.2cm}@{}}
\toprule
\textbf{Nghiên cứu} & \textbf{Năm} & \textbf{Phương pháp chính} & \textbf{Điểm mạnh} & \textbf{Khoảng trống / Hạn chế} \\
\midrule
Marciano \& Adelman (PACE) \cite{marciano2017}
  & 2017
  & So sánh sáu thuật toán ML (SVM, MLP, LR, kNN, DT) trên đặc trưng thống kê STR; lựa chọn đặc trưng bằng KL divergence; phân loại dưới 1 giây
  & Không dùng dữ liệu mô phỏng; ${\sim}97\%$ accuracy trên tập kiểm thử thực; khung so sánh đa thuật toán rõ ràng; thực thi thời gian thực
  & Chỉ bốn lớp NOC (1–4 người); bộ kit hẹp (1.405 profile Identifiler, 20 cá thể); đặc trưng thủ công không mô hình hóa cấu trúc nội locus; khó chuyển sang bộ kit GlobalFiler hay PROVEDIt \\
\addlinespace
Alamoudi et al.\ (TAWSEEM) \cite{alamoudi2022}
  & 2022
  & MLP sâu trên véc tơ phẳng ${\sim}$6.000 chiều bao phủ bốn multiplex STR từ PROVEDIt; đánh giá ba kịch bản tập dữ liệu khác nhau về quy mô
  & $97\%$ accuracy trong kịch bản bốn multiplex với 1–5 người; huấn luyện trên PROVEDIt quy mô lớn; cho thấy MLP sâu vượt random forest trên dữ liệu lớn
  & Véc tơ đầu vào phụ thuộc chặt vào bộ kit và số allele bin; phương pháp phân chia dữ liệu không ngăn rò rỉ mixture identity (không GroupKFold); chưa công bố mã nguồn để tái lập kết quả \\
\addlinespace
Taylor \& Humphries (deepNoC) \cite{taylor2024deepnoc}
  & 2024
  & Mạng sâu DNN huấn luyện trên 100.000 profile STR được mô phỏng; fine-tune thêm bằng profile casework thực của phòng thí nghiệm; đầu ra phụ hỗ trợ giải thích
  & Mở rộng NOC đến 10 người đóng góp; chiến lược transfer learning từ mô phỏng sang thực nghiệm; cho phép thích nghi với điều kiện riêng từng phòng lab
  & Chất lượng mô phỏng quyết định chất lượng tiền huấn luyện; khoảng cách giữa synthetic signal và casework thật chưa được định lượng; mô hình không khai thác cấu trúc tập peak theo locus \\
\bottomrule
\end{tabular}%
}
\end{table}

\subsection{Phương pháp thống kê và quy tắc đếm allele}

Maximum allele count là cận dưới dễ tính cho NOC, nhưng hiệu năng suy giảm khi các contributor chia sẻ allele hoặc khi dropout làm allele thật biến mất. Haned và cộng sự đề xuất tính likelihood cho các giả thuyết số contributor bằng cả số lượng, độ lớn và tần suất allele, nhờ đó biến NOC thành bài toán so sánh mô hình xác suất \cite{haned2011}. Alfonse và cộng sự sau đó so sánh NOCIt, maximum likelihood estimator và maximum allele count trên 100 mixture phức tạp, tập trung vào giới hạn low template \cite{alfonse2017}. NOCIt được kiểm chứng quy mô lớn trên 815 profile GlobalFiler gồm 100 mẫu một người, 193 mẫu hai người, 170 mẫu ba người, 186 mẫu bốn người và 166 mẫu năm người trong một nghiên cứu validation \cite{grgicak2020}. Điểm mạnh của các phương pháp xác suất là có thể liên hệ trực tiếp với likelihood ratio, nhưng chi phí tính toán, giả định mô hình và nhu cầu cấu hình chuyên gia khiến chúng không phải lúc nào cũng phù hợp cho sàng lọc nhanh.

EuroForMix sử dụng mô hình liên tục để đánh giá hồ sơ STR có artefact và đã trở thành một tham chiếu quan trọng trong probabilistic genotyping nguồn mở \cite{bleka2016}. STRmix cũng được developmental validation cho diễn giải mixture và được sử dụng rộng rãi trong thực hành \cite{bright2016}. Các hệ thống này giải quyết bài toán đánh giá chứng cứ toàn diện hơn một classifier NOC, nhưng NOC thường vẫn phải được đặt trước hoặc đánh giá qua nhiều giả thuyết. Do đó, một mô hình dự đoán NOC nhanh không thay thế probabilistic genotyping mà đóng vai trò đề xuất giả thuyết và kiểm tra tính hợp lý trước khi chạy mô hình xác suất.

\subsection{Machine learning cổ điển cho Number of Contributors}

PACE của Marciano và Adelman là một trong những hệ thống machine learning sớm cho NOC. Công trình sử dụng 1.405 profile không mô phỏng tạo từ 20 cá thể, một đến bốn contributor, 15 locus Identifiler, và đánh giá sáu thuật toán gồm nearest neighbor, decision tree, logistic regression, multilayer logistic regression, multilayer perceptron và support vector machine. Kullback Leibler divergence được dùng để loại đặc trưng có divergence dưới 0,01. Hệ thống báo cáo khoảng 98\% accuracy ở training và 97\% ở testing, đồng thời thực hiện phân loại trong thời gian dưới một giây \cite{marciano2017,alamoudi2022}. Hạn chế chính là dữ liệu và bộ kit hẹp hơn PROVEDIt, trong khi feature engineering được thiết kế riêng cho pipeline.

Benschop và cộng sự xây dựng tập 590 PowerPlex Fusion 6C profile từ 1.174 donor, bao phủ một đến năm contributor và các mức allele sharing, dropout, mixture proportion cùng lượng DNA khác nhau. Sau 50 lần huấn luyện và kiểm thử mười classifier, random forest với 19 đặc trưng được chọn làm RFC19 \cite{benschop2019,alamoudi2022}. Kruijver và cộng sự phát triển decision tree minh bạch trên 766 mixture PROVEDIt và cho thấy hiệu năng phụ thuộc mạnh vào bộ lọc stutter cùng artefact; khoảng accuracy được TAWSEEM tổng hợp cho thiết lập này là 0,77 đến 0,85 \cite{kruijver2021,alamoudi2022}. Ưu điểm của decision tree là tốc độ và khả năng giải thích trực tiếp, nhưng việc tóm tắt hồ sơ thành số ít đặc trưng có thể làm mất cấu trúc peak.

\subsubsection{XGBoost}

XGBoost của Chen và Guestrin là hệ thống tree boosting được thiết kế để tối ưu hiệu suất và khả năng mở rộng thông qua histogram-based splitting và xấp xỉ gradient bậc hai \cite{chen2016}. Mô hình xây dựng tập $T$ cây quyết định $f_1, \ldots, f_T$ theo cơ chế cộng dồn:
\begin{equation}
\hat{y}_i = \sum_{t=1}^{T} f_t(\bm{x}_i),
\end{equation}
trong đó mỗi cây $f_t : \mathbb{R}^d \to \mathbb{R}$ ánh xạ véc tơ đặc trưng sang giá trị lá. Khác với gradient boosting cổ điển, XGBoost thêm số hạng chuẩn hóa vào hàm mục tiêu để kiểm soát độ phức tạp cây:
\begin{equation}
\mathcal{L}^{(t)} = \sum_{i=1}^{n} l\!\left(y_i,\; \hat{y}_i^{(t-1)} + f_t(\bm{x}_i)\right) + \Omega(f_t),
\qquad \Omega(f) = \gamma T + \tfrac{1}{2}\lambda \|w\|^2,
\end{equation}
với $T$ là số lá, $w$ là véc tơ trọng số lá, $\gamma$ phạt số lá tạo độ thưa cây và $\lambda$ phạt chuẩn L2. Tại mỗi bước boosting, hàm mất mát được xấp xỉ Taylor bậc hai theo dự đoán hiện tại $\hat{y}_i^{(t-1)}$; trọng số lá tối ưu tại lá $j$ thu được theo công thức dạng đóng:
\begin{equation}
w_j^* = -\frac{G_j}{H_j + \lambda},
\qquad G_j = \sum_{i \in I_j} g_i, \quad H_j = \sum_{i \in I_j} h_i,
\end{equation}
trong đó $g_i = \partial l/\partial \hat{y}_i^{(t-1)}$ là gradient bậc nhất và $h_i = \partial^2 l/\partial(\hat{y}_i^{(t-1)})^2$ là Hessian, còn $I_j$ là tập chỉ số mẫu thuộc lá $j$. Điểm phân tách tốt nhất được tìm bằng lược đồ histogram: đặc trưng liên tục được rời rạc hóa thành tối đa 256 bin, giúp tìm split trong $O(Bd)$ với $B$ số bin thay vì $O(nd)$ với $n$ số mẫu. Trong nghiên cứu này, XGBoost được cấu hình với 300 cây, độ sâu tối đa sáu và learning rate 0,05, đạt Macro F1 bằng 0,9230 và Micro F1 bằng 0,9763 trên GlobalFiler GroupKFold năm phần.

\emph{Điểm mạnh:} XGBoost xử lý dữ liệu thưa và giá trị thiếu tự nhiên qua cơ chế sparsity-aware split finding, phù hợp với hồ sơ STR có nhiều bin allele trống. Chuẩn hóa L2 và kiểm soát số lá giúp giảm overfitting mà không cần tìm kiếm validation phức tạp. Kết quả Macro F1 0,9230 xác nhận dữ liệu phẳng vẫn chứa tín hiệu NOC đáng kể và đặt ra một baseline khó cho mô hình học sâu. \emph{Hạn chế:} Đặc trưng đầu vào phải được phẳng hóa thành véc tơ phụ thuộc bộ kit, trong khi quan hệ peak-locus và tỷ số stutter phải được tóm tắt thủ công; cấu trúc tập hợp tự nhiên của dữ liệu STR không được mô hình hóa trực tiếp.

\subsubsection{TabPFN}

TabPFN của Hollmann và cộng sự tiếp cận tabular classification theo hướng hoàn toàn khác: thay vì tối ưu tham số trên từng tập dữ liệu, một Transformer được tiền huấn luyện trên hàng triệu tập dữ liệu bảng tổng hợp để học xấp xỉ suy luận Bayes \cite{hollmann2023}. Tại thời điểm suy luận, toàn bộ tập huấn luyện $\mathcal{D} = \{(\bm{x}_i, y_i)\}_{i=1}^n$ và mẫu kiểm thử $\bm{x}^*$ được nối thành một chuỗi và đưa vào Transformer như ngữ cảnh in-context; mô hình xuất xác suất hậu nghiệm xấp xỉ mà không cần bất kỳ training loop nào:
\begin{equation}
p(y \mid \bm{x}^*, \mathcal{D}) \approx q_\theta(y \mid \bm{x}^*, \mathcal{D}),
\end{equation}
trong đó $q_\theta$ là Transformer với tham số $\theta$ cố định sau tiền huấn luyện. Prior được xây dựng bằng cách lấy mẫu từ không gian siêu tham số của Bayesian neural networks ngẫu nhiên và sinh dữ liệu bảng tổng hợp tương ứng, giúp hậu nghiệm xấp xỉ học được biên rộng và tổng quát hóa tốt trên nhiều dạng phân phối đặc trưng. Cơ chế in-context cho phép toàn bộ tập huấn luyện đóng vai trò ngữ cảnh, thay thế cho gradient descent truyền thống bằng một lần forward pass duy nhất.

\emph{Điểm mạnh:} TabPFN đưa ra dự đoán trong chưa một giây mà không cần điều chỉnh siêu tham số; hiệu suất trên tập nhỏ thường cạnh tranh với XGBoost được tinh chỉnh kỹ. \emph{Hạn chế:} Giới hạn kích thước ngữ cảnh của TabPFN gốc khiến tập dữ liệu phải được lấy mẫu con khi vượt quá vài trăm đến vài nghìn mẫu; phiên bản nghiên cứu sử dụng tập con GlobalFiler để phù hợp với giới hạn này. Ngoài ra, mô hình không có cơ chế mô hình hóa cấu trúc peak-locus hay quan hệ stutter một cách tường minh, hạn chế khả năng tận dụng đặc trưng sinh học của hồ sơ STR.

\subsection{Deep learning cho Number of Contributors}

TAWSEEM là một MLP sâu cho NOC được huấn luyện trên PROVEDIt. Tác giả mô tả PROVEDIt có hơn 25.000 multiplex STR profile từ một đến năm người và thử ba kịch bản gồm 780 profile một multiplex, 5.700 profile ba multiplex và 6.000 profile bốn multiplex. Kịch bản lớn nhất sử dụng khoảng 6.000 đặc trưng và đạt 97\% accuracy; tác giả cũng báo cáo 97\% accuracy cho mixture năm contributor \cite{alamoudi2022}. Đây là kết quả nổi bật, nhưng đầu vào phẳng phụ thuộc mạnh vào multiplex và khó chuyển trực tiếp sang kit có số locus khác.

deepNoC giải quyết hạn chế dữ liệu bằng mô phỏng electrophoretic signal. Công trình sinh 100.000 profile gán nhãn trước, huấn luyện mạng sâu cho nhãn từ một đến mười contributor và báo cáo 89\% accuracy. Tác giả cho thấy có thể fine tune với vài trăm profile để đạt accuracy tương tự trong một phòng thí nghiệm cụ thể, đồng thời bổ sung các đầu ra phụ nhằm hỗ trợ giải thích \cite{taylor2024deepnoc}. So với TAWSEEM, deepNoC mở rộng miền nhãn nhưng phụ thuộc chất lượng của simulator; khoảng cách giữa synthetic signal và casework thật vẫn cần được định lượng.

\subsubsection{CNN một chiều}

Mạng tích chập một chiều là baseline học sâu phổ biến cho bài toán phân loại trên véc tơ đặc trưng có cấu trúc chuỗi. Với hồ sơ STR đã phẳng hóa theo cột allele thành véc tơ $\bm{x} \in \mathbb{R}^d$, phép tích chập một chiều tại lớp thứ $l$ với kernel kích thước $k$ và bộ lọc thứ $j$ là:
\begin{equation}
h_{lj}[n] = \operatorname{ReLU}\!\left(\operatorname{BN}\!\left(\sum_{m=0}^{k-1} W_{lj}[m]\, h_{l-1}[n-m] + b_{lj}\right)\right),
\end{equation}
trong đó $\operatorname{BN}$ là Batch Normalization chuẩn hóa kích hoạt trên mini-batch để ổn định quá trình huấn luyện. Parameter sharing trong tích chập cho phép mỗi kernel phát hiện một mẫu đặc trưng cục bộ ở mọi vị trí của chuỗi, giảm đáng kể số tham số so với MLP toàn kết nối. Kiến trúc CNN một chiều của nghiên cứu gồm ba khối tích chập với số kênh lần lượt là 64, 128 và 256; mỗi khối theo thứ tự \texttt{Conv1d}(kernel=3, padding=1) -- \texttt{BatchNorm1d} -- \texttt{ReLU} -- Pooling. Hai khối đầu dùng \texttt{MaxPool1d(2)} giảm nửa chiều dài chuỗi; khối ba dùng \texttt{AdaptiveAvgPool1d(8)} ép chiều dài về đúng tám vị trí bất kể kích thước đầu vào, tạo biểu diễn $256\times8=2048$ chiều sau \texttt{Flatten}. MLP phân loại: Linear($2048\to256$)--ReLU--Dropout(0{,}3) -- Linear($256\to128$)--ReLU--Dropout(0{,}2) -- Linear($128\to5$); softmax cho xác suất năm lớp NOC. Mô hình tối ưu bằng Adam lr=$10^{-3}$, weight decay $10^{-4}$, dừng sớm patience 20 epoch theo Macro F1 validation.

\emph{Điểm mạnh:} Tích chập học đặc trưng cục bộ với parameter sharing và tính bất biến dịch chuyển cục bộ; Batch Normalization giúp huấn luyện ổn định kể cả khi learning rate lớn. \emph{Hạn chế:} Véc tơ phẳng hóa không có cấu trúc tường minh về locus; allele của cùng locus nhưng nằm xa nhau trong cách sắp xếp bộ kit có thể bị kernel bỏ qua, trong khi quan hệ stutter và cân bằng height trong locus đòi hỏi đặc trưng thủ công. Trên GlobalFiler GroupKFold năm phần, CNN một chiều đạt Macro F1 bằng 0,8854, thấp hơn DS-ST gần 0,059 điểm tuyệt đối.

\subsubsection{Forensic Transformer}

Forensic Transformer mô hình hóa hồ sơ STR ở mức locus thay vì mức allele, xây dựng trực tiếp trên kiến trúc Transformer \cite{vaswani2017}. Hai mươi bốn locus được xem là hai mươi bốn token độc lập. Một đặc điểm quan trọng trong hiện thực là mỗi locus $m$ có số cột allele $d_m$ khác nhau (ví dụ SE33 có nhiều allele hơn TH01), nên mỗi locus được chiếu bởi một lớp tuyến tính riêng: $\mathrm{Linear}(d_m \to 64)$ với $d_m = \max(|\mathcal{A}_m|, 1)$. Kết quả là 24 token $\bm{Z} \in \mathbb{R}^{24 \times 64}$ không cần positional embedding bổ sung vì thứ tự locus cố định theo bảng định nghĩa STR. Chuỗi 24 locus token đi qua ba \texttt{TransformerEncoderLayer}, mỗi lớp thực hiện multi-head self-attention với bốn head:
\begin{equation}
\operatorname{head}_r = \operatorname{softmax}\!\left(\frac{Q_r K_r^\top}{\sqrt{d_k}}\right) V_r, \quad d_k = \frac{d_{\mathrm{locus}}}{h} = 16,
\end{equation}
\begin{equation}
\operatorname{MHA}(Q,K,V) = \operatorname{Concat}(\operatorname{head}_1, \ldots, \operatorname{head}_h)\, W^O,
\end{equation}
theo sau bởi add-and-norm residual và feed-forward network kích thước 256 với GELU, dropout 0{,}1, batch\_first=True. Sau ba lớp, \texttt{mean} trên chiều locus tổng hợp 24 token thành biểu diễn hồ sơ 64 chiều; đầu phân loại gồm \texttt{LayerNorm}(64)$\to$\texttt{Linear}(64$\to$5) và chiếu xuống năm logit NOC. Mô hình tối ưu bằng Adam lr=$5\times10^{-4}$, weight decay $10^{-5}$, dừng sớm patience 20 theo Macro F1.

\emph{Điểm mạnh:} Self-attention theo cặp locus cho phép mô hình học quan hệ toàn hồ sơ không phụ thuộc khoảng cách vị trí; locus có stutter phức tạp có thể được hiệu chỉnh bởi attention từ các locus sạch khác. \emph{Hạn chế:} Việc tóm tắt toàn bộ peak của một locus thành một token duy nhất trước bước attention có thể làm mất thông tin về contributor thiểu số có peak yếu; peak thiểu số bị pha loãng trong phép cộng tổng hợp tại locus, khiến mô hình khó phân biệt NOC khi một contributor có tỷ lệ hỗn hợp thấp. Trên GlobalFiler GroupKFold năm phần, Forensic Transformer đạt Macro F1 bằng 0,8846, xấp xỉ CNN một chiều và thấp hơn XGBoost 0,038 điểm.

Veldhuis và cộng sự tập trung vào explainability. Thay vì chỉ đưa ra feature importance, họ xây dựng counterfactual explanation thực tế cho một dự đoán NOC, nghĩa là tạo thay đổi có thể xảy ra trong hồ sơ DNA và đủ làm mô hình đổi nhãn \cite{veldhuis2022}. Hướng này đặc biệt quan trọng vì một classifier đạt accuracy cao nhưng không giải thích được yếu tố peak hoặc locus nào dẫn đến quyết định sẽ khó được kiểm chứng trong môi trường pháp y.

DS-ST khác các MLP nói trên ở chỗ đầu vào không bị ép thành một véc tơ duy nhất. Mỗi peak còn tồn tại như một token độc lập sau bước enrichment; hai ISAB++ học quan hệ giữa peak trên toàn hồ sơ với chi phí \(O(NM)\), còn PMA học phép tổng hợp thành véc tơ mixture. Quan hệ stutter không được cài thành một attention bias giả định trước mà được đưa vào trường \(SR=h/h_{\mathrm{parent}}\), nhờ đó mô hình có thể quyết định cách sử dụng tỷ số này cùng các đặc trưng peak khác.

\subsection{Deep learning cho nhận dạng contributor và allele calling}

Phan và cộng sự đề xuất pipeline học sâu dự đoán cá thể trong mixture sequencing, sử dụng 27 STR và 94 SNP trên dữ liệu từ sáu cá thể \cite{phan2021}. Bài toán này gần contributor identification hơn NoC, nhưng dữ liệu massively parallel sequencing giàu thông tin sequence hơn capillary electrophoresis. Vì vậy, hiệu năng của công trình không thể so sánh trực tiếp với mô hình xử lý peak height.

DNANet chuyển allele calling thành segmentation trên ma trận electropherogram kích thước năm dye lane nhân 4.096 scan point. Tác giả khảo sát 90 cấu hình tạo bởi ba độ sâu mạng, hai lựa chọn kernel và năm số filter khởi đầu. Phiên bản xuất bản báo cáo F1 bằng 0,971 trên allele do analyst gán nhãn trong case data chưa thấy, 0,982 trên research data và 0,962 khi so với allele donor thật, bằng F1 của annotation analyst trên tiêu chí này \cite{dewit2025}. DNANet không dự đoán NOC, nhưng cho thấy việc giữ cấu trúc tín hiệu thay vì chỉ dùng feature tổng hợp có giá trị thực nghiệm.

Khung đề xuất trong nghiên cứu này nằm sau bước allele calling. Đầu vào là các peak đã được trích xuất và chuẩn hóa, trong khi DNANet có thể trở thành module phía trước để tự động tạo peak đáng tin cậy. Sự kết hợp tiềm năng giữa DNANet và DS-ST hình thành pipeline từ electropherogram thô đến allele, NOC và contributor, nhưng việc tích hợp đầy đủ nằm ngoài phạm vi thực nghiệm hiện tại.

\subsection{Mạng học trên tập hợp và attention cảm ứng}

Transformer chuẩn sử dụng scaled dot product attention
\begin{equation}
\operatorname{Attention}(Q,K,V)=
\operatorname{softmax}\left(\frac{QK^\top}{\sqrt{d_k}}\right)V,
\end{equation}
và loại bỏ recurrence trong kiến trúc gốc \cite{vaswani2017}. Khi không sử dụng positional encoding, self attention có tính equivariant theo hoán vị, phù hợp với tập peak. Set Transformer xây dựng Multihead Attention Block, Induced Set Attention Block và Pooling by Multihead Attention để học hàm trên tập \cite{lee2019}. SetNorm bổ sung chuẩn hóa dành riêng cho tập và clean path residual, giúp các mạng set sâu ổn định hơn \cite{zhang2023setnorm}.

Deep Sets tạo biểu diễn bất biến bằng phép cộng
\begin{equation}
\bm{z}=\rho\left(\sum_{i=1}^{n}\phi(\bm{x}_i)\right).
\end{equation}
Trong mô hình của đồ án, PMA thay phép cộng cố định bằng một seed vector học được. Seed này truy vấn toàn bộ tập peak đã mã hóa và tạo biểu diễn hồ sơ 128 chiều. Nhờ không sử dụng positional encoding theo thứ tự cột, phép mã hóa giữ tính equivariant theo hoán vị ở mức token, còn PMA tạo đầu ra invariant theo hoán vị. Đây là điểm khác biệt quan trọng với CNN một chiều và với Transformer chia cột locus trong baseline.

\subsection{Slot attention và optimal transport cho contributor identification}

Slot Attention tạo một tập véc tơ slot và cập nhật chúng qua nhiều vòng attention cạnh tranh \cite{locatello2020slot}. Trong thị giác, mỗi slot thường biểu diễn một vật thể; trong nghiên cứu này, mỗi slot được gắn với một donor tham chiếu. Conditional Slot Attention sử dụng grounded slot dictionary để khởi tạo slot có điều kiện thay vì lấy mẫu hoàn toàn ngẫu nhiên \cite{ma2024cosa}. Tương tự, genotype của từng donor được mã hóa thành slot khởi tạo, giúp donor slot bắt đầu từ allele tham chiếu của chính donor đó.

Guided Slot Attention sử dụng thông tin hướng dẫn để cải thiện phân tách foreground và background trong video \cite{lee2024gsa}. Khi điều chỉnh sang DNA, logits gán peak cho donor đóng vai trò hướng dẫn. MESH liên hệ Slot Attention với optimal transport và tối thiểu hóa entropy của Sinkhorn để phá hòa giữa các slot tương tự \cite{zhang2023mesh}. Điều này phù hợp với trường hợp hai donor chia sẻ nhiều allele. Adaptive Slot Attention giải quyết số slot biến đổi bằng cơ chế chọn slot rời rạc và masked decoder \cite{fan2024adaslot}; trong bối cảnh DNA, gate tồn tại của slot là tín hiệu trực tiếp cho việc donor có mặt hay không.

\section{Mô hình đề xuất}

\subsection{Tập dữ liệu}

\subsubsection{Nguồn dữ liệu PROVEDIt}

PROVEDIt (Project Research Openness for Validation with Empirical Data) là một tập dữ liệu STR pháp y quy mô lớn được công bố nhằm hỗ trợ phát triển, so sánh và kiểm định các phương pháp diễn giải hỗn hợp DNA; bài báo gốc mô tả tập dữ liệu gồm cả hồ sơ nguồn đơn và hồ sơ hỗn hợp \cite{provedit2018}, còn trang dữ liệu mở của Đại học Rutgers hiện công bố hơn 27.000 hỗn hợp DNA pháp y \cite{proveditweb}. Về loại dữ liệu, đây là tín hiệu điện di mao quản (capillary electrophoresis): mỗi hồ sơ là một tập peak huỳnh quang, trong đó mỗi peak gắn với một locus STR, một giá trị allele và một chiều cao tính bằng đơn vị RFU; tập gốc được tạo trong hơn bốn năm qua 144 điều kiện phòng thí nghiệm, trải từ một đến năm người đóng góp và có cả ba bản raw, filtered, unfiltered \cite{alamoudi2022}. Về định dạng, mỗi hồ sơ ở dạng phân phối là một bảng cột chiều cao allele theo từng locus; nghiên cứu chuyển mỗi hồ sơ thành một tập token peak ở dạng tensor \((N,160,8)\), nghĩa là tối đa \(N=160\) peak và tám trường số cho mỗi peak, kèm một mặt nạ đánh dấu peak hợp lệ.

Nghiên cứu sử dụng ba nhánh dữ liệu thực được tạo từ các bộ kit GlobalFiler, Identifiler Plus và PowerPlex 16. GlobalFiler gồm 10.124 hồ sơ và 439 cột tính cả định danh nhóm cùng nhãn, Identifiler Plus gồm 3.160 hồ sơ và 281 cột, còn PowerPlex 16 gồm 1.022 hồ sơ và 214 cột (Bảng~\ref{tab:data_stats}). Sự khác biệt số cột phản ánh số locus và allele bin khác nhau giữa các kit, là lý do kiến trúc không nên phụ thuộc một véc tơ đầu vào có chiều cố định.

\begin{table}[H]
\centering
\caption{Thống kê các tập dữ liệu đã xây dựng từ PROVEDIt.}
\label{tab:data_stats}
\small
\begin{tabular}{lrrrrrrr}
\toprule
\textbf{Bộ kit} & \textbf{Hồ sơ} & \textbf{Nhóm} & \textbf{NOC 1} & \textbf{NOC 2} &
\textbf{NOC 3} & \textbf{NOC 4} & \textbf{NOC 5} \\
\midrule
GlobalFiler & 10.124 & 3.418 & 8.119 & 526 & 484 & 527 & 468 \\
Identifiler Plus & 3.160 & 1.084 & 2.231 & 365 & 209 & 146 & 209 \\
PowerPlex 16 & 1.022 & 435 & 794 & 56 & 52 & 60 & 60 \\
\bottomrule
\end{tabular}
\end{table}

Phân phối GlobalFiler có 80,20\% hồ sơ thuộc lớp một contributor, trong khi lớp hiếm nhất là năm contributor chỉ chiếm 4,62\%. Nếu tối ưu accuracy thông thường, mô hình có thể đạt điểm cao bằng cách ưu tiên NOC một nhưng thất bại ở mixture phức tạp. Vì vậy, Macro F1 được sử dụng làm metric lựa chọn chính, Focal Loss giảm trọng số mẫu dễ và oversampling được thực hiện chỉ trong tập huấn luyện.

\subsubsection{Dữ liệu in-silico cho nhận dạng contributor}

Nhận dạng contributor cần biết chính xác tập người đóng góp, và để giám sát attribution theo peak còn cần biết mỗi peak thuộc về donor nào — thông tin không tồn tại cho hỗn hợp thực. Vì vậy, dữ liệu huấn luyện DS-ST là các hỗn hợp bán-tổng hợp được tạo bằng cách chồng phủ (overlay) các hồ sơ \emph{nguồn đơn thực} (một người, NOC \(=1\)) của \(C=45\) donor trong panel tham chiếu, theo mô hình chiều cao peak trong giám định \cite{bleka2016}.

Cụ thể, để tạo một hỗn hợp \(k\) người, trước hết lấy mẫu một véc tơ tỷ lệ đóng góp \(\bm{\phi}=(\phi_1,\dots,\phi_k)\) với \(\sum_c\phi_c=1\) và một mức khuôn tổng \(T\). Với mỗi donor \(c\) được chọn, lấy ngẫu nhiên một hồ sơ nguồn đơn thực của \(c\), đổi về chiều cao RFU \(\bm{h}=\exp(\cdot)-1\), chuẩn hóa về chiều cao tương đối \(\bm{h}\leftarrow\bm{h}/\textstyle\sum_b h_b\) vì chiều cao tuyệt đối phụ thuộc lượng khuôn, nhân thêm một nhiễu gamma theo từng peak, rồi cộng đóng góp \(\phi_c\,T\,\bm{h}\) vào hỗn hợp; các allele dùng chung tự cộng dồn:
\begin{equation}
\bm{m}=\sum_{c=1}^{k}\phi_c\,T\,\bm{h}^{(c)},\qquad \bm{h}^{(c)}=\frac{\exp(\bm{ss}^{(c)})-1}{\sum_b(\exp(\bm{ss}^{(c)})-1)_b}\odot \bm{\varepsilon}^{(c)},
\end{equation}
với \(\bm{ss}^{(c)}\) là hồ sơ nguồn đơn thực và \(\bm{\varepsilon}^{(c)}\sim\mathrm{Gamma}\) là nhiễu chiều cao. Vì phép chuẩn hóa làm mất một phần stutter của thành phần thiểu số, một số hạng back-stutter được thêm trên hỗn hợp đã cộng: tại allele \(a-1\) cộng \(SR(a)\) lần chiều cao peak cha, với \(SR\) tuyến tính theo giá trị allele và nhiễu log-chuẩn, được hiệu chỉnh để tỷ lệ stutter cuối khớp dữ liệu thực. Cuối cùng, mọi peak dưới ngưỡng phát hiện \(AT\) bị đưa về không, tái hiện hiện tượng dropout của contributor thiểu số. Toàn bộ quá trình dùng một seed cố định; nhãn donor, tỷ lệ \(\bm{\phi}\) và donor trội tại mỗi peak (nhãn provenance) được ghi lại chính xác nhờ biết đóng góp của từng nguồn. Mỗi mẫu là một tập tối đa \(160\) peak ở định dạng tensor \((N,160,8)\) kèm mặt nạ hợp lệ, đồng nhất với biểu diễn của nhánh đếm.

\subsubsection{Phân chia tập huấn luyện, kiểm thử và thử nghiệm}

Phần thử nghiệm (test) gồm các hỗn hợp \emph{thực} của PROVEDIt và đóng vai trò tập đánh giá vàng. Điểm mấu chốt để chống rò rỉ là các \emph{tổ hợp người đóng góp} xuất hiện trong phần thử nghiệm được loại trừ khỏi quá trình tổng hợp in-silico, nên không một tổ hợp donor nào của tập thử nghiệm từng được mô hình nhìn thấy khi huấn luyện. Phần kiểm thử (validation) gồm các hỗn hợp thực còn lại không nằm trong phần thử nghiệm và được dùng để lựa chọn mô hình. Phần huấn luyện (training) gồm toàn bộ hỗn hợp in-silico nói trên, với các tổ hợp của tập thử nghiệm đã bị loại. Cách chia này đo trực tiếp khả năng tổng quát hóa sang tổ hợp người đóng góp mới: mô hình học trên hỗn hợp bán-tổng hợp nhưng được đánh giá trên hỗn hợp thực với tổ hợp tách biệt hoàn toàn.

Về quy mô, phần huấn luyện gồm 55.247 hỗn hợp in-silico, phần kiểm thử gồm 1.366 và phần thử nghiệm gồm 3.582 hỗn hợp thực; ngoài ra một phân vùng open-set gồm 1.366 hỗn hợp có ít nhất một contributor ngoài panel được dùng riêng để huấn luyện và đánh giá đầu từ chối. Phần huấn luyện được lấy mẫu nghiêng về các lớp NOC cao, phân bố xấp xỉ 5.247, 7.266, 10.625, 17.709 và 14.400 mẫu cho một đến năm người, nhằm cung cấp đủ tín hiệu cho các hỗn hợp khó; trong khi phần kiểm thử và phần thử nghiệm thực giữ phân bố lệch mạnh về nguồn đơn đúng như casework, với phần thử nghiệm gồm lần lượt 2.249, 334, 385, 242 và 372 mẫu cho một đến năm người. Việc lựa chọn mô hình căn cứ trên recall oracle trung bình theo NOC đo trên phần kiểm thử, còn mọi chỉ số công bố đều được tính trên phần thử nghiệm thực tách biệt.

Riêng cho bài toán Number of Contributors, mô hình đề xuất và các baseline được so sánh trên dữ liệu GlobalFiler thực với GroupKFold năm phần theo định danh mixture: \texttt{group\_id} được dựng bằng cách loại thành phần thời gian tiêm khỏi tên mẫu, nhờ đó các biến thể gần trùng của cùng một mixture không rơi đồng thời vào phần huấn luyện và phần kiểm thử.

\subsection{Biểu diễn đầu vào}

DS-ST biểu diễn mỗi hồ sơ như một tập peak phẳng thay vì một tensor cố định theo locus. Cụ thể, hồ sơ là một tập tối đa \(N=160\) token peak, mỗi token là một véc tơ tám trường, gộp thành tensor
\begin{equation}
\mathcal{X}\in\mathbb{R}^{B\times N\times 8},
\end{equation}
kèm mặt nạ \(\mathcal{M}\in\{0,1\}^{B\times N}\) đánh dấu peak hợp lệ; chỉ số locus nằm ngay trong token nên cấu trúc locus được bảo toàn mà không cần xếp peak vào một lưới locus cố định, và do đó mô hình không phụ thuộc số locus hay số allele bin của từng bộ kit.

Tám trường của token được chắt lọc từ một tập đặc trưng kỹ thuật rộng hơn — gồm định danh locus và allele, log chiều cao peak, tỷ lệ chiều cao so với tổng locus, chiều cao tương đối toàn hồ sơ, thứ hạng peak trong locus, số peak tại locus, khoảng cách tới allele lân cận và các thống kê giúp nhận biết stutter — giữ lại những trường mang nhiều thông tin nhất cho các nhiệm vụ:
\begin{equation}
\bm{x}_{lp}=
\left[l,\ a,\ \log(1+h),\ H_b,\ SR,\ r,\ n_{\mathrm{peak}},\ h_{\mathrm{rel}}\right].
\end{equation}
Trong đó \(H_b=h/\sum_{q\in l}h_q\), \(SR=h/h_{a+1}\) khi peak cha tồn tại, \(r\) là rank chuẩn hóa, \(n_{\mathrm{peak}}\) là số peak chuẩn hóa và \(h_{\mathrm{rel}}=h/\max_q h_q\).

\subsection{Tổng quan kiến trúc Donor-Slot Set Transformer}

Mô hình đề xuất, gọi là Donor-Slot Set Transformer (DS-ST), xem mỗi hồ sơ STR là một tập peak bất biến đối với hoán vị và xử lý toàn bộ tập bằng một bộ mã hóa attention duy nhất, thay vì nén từng locus thành một véc tơ trước khi suy luận. Mỗi peak được biểu diễn bằng một token tám trường, được chiếu vào không gian \(d_{\mathrm{model}}=128\), được lọc bỏ nếu không khớp bất kỳ kiểu gen tham chiếu nào, rồi đưa qua hai khối Induced Set Attention cải tiến để thu biểu diễn ngữ cảnh \(\mathbf{H}\in\mathbb{R}^{B\times N\times128}\) cho mọi peak. Từ \(\mathbf{H}\), một bộ giải mã khe có điều kiện kiểu gen sinh một xác suất cho mỗi donor trong panel tham chiếu \(C=45\) (nhiệm vụ chính CLS), và một biểu diễn hồ sơ gộp từ nhánh không lọc cho điểm từ chối open-set (nhiệm vụ chính reject). Số người đóng góp không phải mục tiêu chính mà được suy ra từ chính phân bố xác suất donor ở bước hậu xử lý, cùng với ước lượng tỷ lệ hỗn hợp hỗ trợ huấn luyện. Như vậy một bộ mã hóa duy nhất phục vụ hai nhiệm vụ chính và còn cho phép đọc ra số người như một đại lượng phái sinh.

\begin{figure}[H]
\centering
\resizebox{\textwidth}{!}{
\begin{tikzpicture}[node distance=0.5cm and 0.5cm]
  \node[block,fill=peakblue!12] (tok) {Token tám trường\\\((B,N,8)\)};
  \node[block,fill=peakblue!12,right=of tok] (emb) {Mã hóa token\\LocusEmbed \(+\) PLR\\\(+\) feas\_filter};
  \node[block,fill=locusgreen!12,right=of emb] (enc) {ISAB++ \(\times2\)\\set\_of\_set\\\(\to\mathbf{H}\,(B,N,128)\)};
  \node[block,fill=profileorange!14,right=of enc] (dec) {Bộ giải mã khe\\CoSA \(\to\) GSANet\\\(\to\) MESH \(\to\) AdaSlot};
  \node[smallblock,right=of dec] (cls) {\(\sigma(\bm{\ell})\)\\\(P(\mathrm{donor}_c\mid X)\)};
  \node[block,fill=profileorange!10,above=0.7cm of dec] (geno) {Panel kiểu gen\\\(C{=}45\) donor};
  \node[block,fill=headpurple!13,below=0.7cm of enc] (pma) {PMA};
  \node[smallblock,below=0.7cm of dec] (aux) {\(\hat{\bm{\pi}}\) mixture\\reject};
  \node[block,fill=headpurple!13,below=0.7cm of cls] (decode) {phi-rerank\\\(+\) đếm RF};
  \node[smallblock,right=of decode] (out) {Tập contributor\\\(+\) NOC \(k\)};
  \draw[line] (tok)--(emb); \draw[line] (emb)--(enc); \draw[line] (enc)--(dec); \draw[line] (dec)--(cls);
  \draw[line] (geno)--(dec); \draw[line] (enc)--(pma); \draw[line] (pma)--(aux);
  \draw[line] (cls)--(decode); \draw[line] (decode)--(out);
\end{tikzpicture}}
\caption{Kiến trúc Donor-Slot Set Transformer. Một bộ mã hóa tập duy nhất chuyển tập peak thành biểu diễn ngữ cảnh \(\mathbf{H}\); bộ giải mã khe có điều kiện kiểu gen trả về xác suất từng donor, còn nhánh gộp cung cấp tỷ lệ hỗn hợp và điểm từ chối. Số người đóng góp được suy ra ở bước giải mã hậu nghiệm.}
\label{fig:overall}
\end{figure}

Hình~\ref{fig:overall} mô tả luồng top-down. Điểm khác biệt cốt lõi so với các kiến trúc theo locus là DS-ST không tổng hợp các peak của một locus thành một token trước attention; mỗi peak vẫn là một phần tử độc lập cho tới khi bộ giải mã phân bổ bằng chứng về donor. Nhờ đó, peak yếu của một contributor thiểu số không bị trung bình hóa cùng các peak cao của contributor chính, vốn là nguyên nhân chính khiến mô hình bỏ sót người đóng góp tỷ lệ thấp. Bộ mã hóa, bộ giải mã khe, các đầu phụ và bước giải mã hậu nghiệm được trình bày lần lượt ở các mục tiếp theo; chính bộ mã hóa này, khi ghép với một đầu đếm năm lớp, tạo thành biến thể chuyên cho bài toán Number of Contributors trong Hình~\ref{fig:settransv4}.

\subsection{Từ Set Transformer gốc đến DS-ST}

Set Transformer xử lý một tập đầu vào bằng attention thay cho phép cộng cố định của Deep Sets \cite{lee2019,zaheer2017}. Đơn vị cơ sở là khối attention đa đầu \(\operatorname{MAB}(X,Y)=\operatorname{LN}\!\big(H+\operatorname{FFN}(H)\big)\) với \(H=\operatorname{LN}\!\big(X+\operatorname{Attn}(X,Y,Y)\big)\); khối induced \(\operatorname{ISAB}(X)=\operatorname{MAB}\!\big(X,\operatorname{MAB}(I,X)\big)\) dùng \(M\) điểm cảm ứng học được \(I\) để hạ chi phí từ \(O(N^2)\) xuống \(O(NM)\); và phép gộp \(\operatorname{PMA}\) dùng một véc tơ truy vấn học được để biến tập thành một véc tơ bất biến hoán vị. Kiến trúc này phù hợp tự nhiên với một hồ sơ STR xem như tập peak, nhưng ở dạng nguyên bản nó bộc lộ bốn hạn chế đối với bài toán nhận dạng người đóng góp, và mỗi hạn chế định hướng một lựa chọn thiết kế của DS-ST.

Thứ nhất, attention softmax mang tính cạnh tranh: trọng số trên các khóa cộng lại bằng một, nên một peak chính rất cao hút phần lớn khối lượng attention và làm mờ peak yếu của một contributor tỷ lệ thấp — chính peak hiếm này lại là bằng chứng quyết định cho người đóng góp thiểu số. DS-ST thay bước thu nhận tại điểm cảm ứng bằng attention sigmoid không cạnh tranh, trong đó mỗi cặp truy vấn–khóa có một cổng độc lập, để peak yếu vẫn được nạp vào biểu diễn bất kể peak chính cao đến đâu. Thứ hai, \(\operatorname{PMA}\) chỉ tạo một véc tơ hồ sơ duy nhất, đủ để phân loại số người nhưng không thể phát ra một quyết định có cấu trúc cho từng donor trong panel \(45\) người, cũng không có chỗ để đưa kiểu gen tham chiếu vào. DS-ST thay đầu gộp bằng một bộ giải mã khe gồm \(45\) khe, mỗi khe được neo vào kiểu gen của một donor, nhờ đó đầu ra có cấu trúc theo donor và tiên nghiệm kiểu gen được tích hợp ngay từ khởi tạo. Thứ ba, bộ mã hóa nguyên bản không có cơ chế chia một peak dùng chung cho nhiều donor cùng mang allele đó; một phép gán softmax đơn buộc chọn một donor thắng tất, gây hiện tượng giải thích lấn. DS-ST đưa định tuyến vận tải tối ưu với chuẩn hóa hai phía vào vòng cập nhật khe, cho phép một peak chia sẻ được phân tỷ lệ cho đúng các donor mang nó. Thứ tư, LayerNorm chuẩn hóa từng token độc lập và làm mất thống kê chiều cao tương đối giữa các peak — vốn mang thông tin tỷ lệ hỗn hợp; DS-ST dùng SetNorm để chuẩn hóa theo thống kê toàn tập \cite{zhang2023setnorm}, kết hợp tách tập peak riêng với peak chia sẻ và bộ lọc khả thi để bảo toàn và làm nổi bật bằng chứng đặc thù donor. Cuối cùng, vì số người đóng góp thay đổi giữa các mẫu, một cổng tồn tại rời rạc cho mỗi khe quyết định donor nào thực sự hiện diện, và một đầu từ chối xử lý trường hợp có người ngoài panel. Các mục sau trình bày chi tiết từng thành phần theo đúng trình tự xử lý.

\subsection{Bộ mã hóa tập và biến thể đếm số người đóng góp}

Trước khi mô tả bộ giải mã khe, mục này trình bày bộ mã hóa tập của DS-ST, vì cùng bộ mã hóa đó cũng trực tiếp giải quyết bài toán Number of Contributors. Khi tách riêng cho nhiệm vụ đếm, mô hình được ký hiệu là DS-ST: nó dùng đúng cách mã hóa token và hai khối ISAB++ của DS-ST, nhưng thay bộ giải mã khe bằng một đầu đếm năm lớp gọn nhẹ. Biến thể này được khởi tạo từ checkpoint của chính nhiệm vụ nhận dạng contributor (panel \(45\) lớp), do đó kết quả Number of Contributors báo cáo về sau là hệ quả trực tiếp của biểu diễn mà DS-ST học được, chứ không phải một mạng tách rời.

\begin{figure}[H]
\centering
\resizebox{\textwidth}{!}{
\begin{tikzpicture}[node distance=0.55cm and 0.55cm]
  \node[block,fill=peakblue!12] (input)
    {Token tám chiều\\$(B,N,8)$};
  \node[block,fill=peakblue!12,right=of input] (embed)
    {LocusEmbed\\$+$ PeriodicPLR\\$\to$ InputProj\\$(B,N,128)$};
  \node[block,fill=locusgreen!14,right=of embed] (mab0)
    {mab0\\Sigmoid-\\MABpp};
  \node[block,fill=locusgreen!12,right=of mab0] (mab1)
    {mab1\\MABpp\\softmax};
  \node[block,fill=profileorange!14,right=of mab1] (pma)
    {PMA\\$k{=}1$\\$\bm{z}_{\mathrm{mix}}\ (B,128)$};
  \node[block,fill=headpurple!13,right=of pma] (head)
    {NOC Head\\$128\!\to\!64\!\to\!5$\\Focal Loss};
  \node[smallblock,right=of head] (out)
    {$P(y\!=\!1..5)$};
  \node[groupbox,fit=(mab0)(mab1)] (isab) {};
  \node[above=2pt of isab,font=\tiny\bfseries] {ISABpp $\times$ 2,\ SetNorm,\ clean-path};
  \draw[line] (input)--(embed);
  \draw[line] (embed)--(mab0);
  \draw[line] (mab0)--(mab1);
  \draw[line] (mab1.east)--(pma.west);
  \draw[line] (pma)--(head);
  \draw[line] (head)--(out);
\end{tikzpicture}}
\caption{Kiến trúc DS-ST. Token tám chiều mỗi peak được mã hóa song song qua LocusEmbed và PeriodicPLR rồi chiếu vào $d_{\mathrm{model}}=128$. Hai khối ISAB++ với SigmoidMABpp tại mab0 học quan hệ giữa mọi peak qua 32 inducing point. PMA tổng hợp thành $\bm{z}_{\mathrm{mix}}$ và đầu NOC chiếu xuống năm lớp.}
\label{fig:settransv4}
\end{figure}

Mỗi peak được biểu diễn bằng token tám chiều
\begin{equation}
\bm{x}_p = \bigl[l,\;a,\;\log(1+h),\;H_b,\;SR,\;r,\;n_l/12,\;h/h_{\max}\bigr],
\end{equation}
trong đó $l$ là chỉ số locus ($0\ldots23$), $a$ là giá trị allele, $h$ là chiều cao peak RFU, $H_b=h/\sum_{q\in l}h_q$ là tỷ lệ chiều cao trong locus, $SR=h/(h_{a+1}+\varepsilon)$ là tỷ số với chiều cao peak stutter cha tại allele $a+1$ (bằng không nếu peak cha không tồn tại), $r=\mathrm{rank}_{\downarrow}/(n_l-1)$ là thứ hạng giảm dần chuẩn hóa trong locus (bằng không cho peak cao nhất), $n_l/12$ là số peak chuẩn hóa và $h/h_{\max}$ là chiều cao tương đối so với peak cao nhất toàn hồ sơ. Bảy trường liên tục được chuẩn hóa z-score theo thống kê tập huấn luyện trước khi đưa vào PeriodicPLR.

Chỉ số locus $l$ được mã hóa bằng $\texttt{Embedding}(25,16,\text{padding\_idx}=24)$ cho ra 16 chiều. Bảy trường liên tục sau khi chuẩn hóa z-score được đưa qua PeriodicPLR ($n_{\mathrm{freq}}=8$, $d_{\mathrm{num}}=8$, $\sigma_{\mathrm{init}}=0{,}3$): mỗi đặc trưng $x_j$ sinh $[\sin(2\pi c_{jk}x_j), \cos(2\pi c_{jk}x_j)]_{k=1}^{n_{\mathrm{freq}}}$, chiếu qua Linear$(2n_{\mathrm{freq}}\to d_{\mathrm{num}})$ rồi ReLU, tạo 8 chiều riêng; bảy đặc trưng ghép thành $7\times8=56$ chiều. Tần số $c_{jk}$ được học từ khởi tạo $\mathcal{N}(0,\sigma^2)$ và đặc trưng hóa chu kỳ tuần hoàn tự nhiên của allele, chiều cao và rank. Hai nhánh embedding được ghép thành $16+56=72$ chiều rồi chiếu vào $d_{\mathrm{model}}=128$ qua $W_{\mathrm{in}}\in\mathbb{R}^{128\times72}$; token của các peak được mask được triệt về không trước khi vào encoder.

Mỗi khối ISAB++ gồm hai bước. Bước mab0 dùng SigmoidMABpp, trong đó $M=32$ inducing point thu nhận tín hiệu peak qua attention không cạnh tranh
\begin{equation}
A_{mc}=\sigma\!\left(\frac{Q_mK_c^\top}{\sqrt{d_h}}+b_{\mathrm{nc}}\right),
\qquad b_{\mathrm{nc}}=-\log n_{\mathrm{valid}},
\end{equation}
với $\sigma$ là hàm sigmoid và $b_{\mathrm{nc}}$ là bias tự động điều chỉnh theo số peak hợp lệ trong mẫu. Khác với softmax, các gate này độc lập nhau, cho phép peak thiểu số được chấp nhận vào inducing point mà không bị triệt tiêu bởi peak chiếm ưu thế. Bước mab1 dùng MABpp softmax chuẩn để mỗi peak truy vấn ngược biểu diễn inducing. SetNorm chuẩn hóa trên toàn tập token và clean-path residual ổn định gradient khi xếp chồng nhiều khối. Sau hai ISAB++, PMA với một seed vector học được tạo $\bm{z}_{\mathrm{mix}}\in\mathbb{R}^{128}$. Đầu NOC áp dụng hai lớp tuyến tính với dropout và ReLU:
\begin{equation}
\hat{\bm{p}} = \operatorname{softmax}\!\bigl(W_2\operatorname{ReLU}(W_1\bm{z}_{\mathrm{mix}})\bigr),
\qquad W_1\in\mathbb{R}^{64\times128},\;W_2\in\mathbb{R}^{5\times64}.
\end{equation}
Quá trình fine tune gồm hai giai đoạn với AdamW ($\text{lr}=3\times10^{-4}$, $\lambda=10^{-4}$) và gradient clipping $\|\nabla\|_{\max}=1{,}0$. Giai đoạn một đóng băng toàn bộ tham số backbone và cập nhật chỉ đầu NOC với learning rate $9\times10^{-4}$ trong năm epoch, ổn định giá trị khởi tạo ngẫu nhiên trước khi mở backbone. Giai đoạn hai mở toàn bộ tham số, áp dụng CosineAnnealingLR suốt 60 epoch với $\eta_{\min}=10^{-6}$ và dừng sớm theo Macro F1 validation với patience bằng mười hai. Batch size bằng 64; trước mỗi fold, tập huấn luyện được oversample để cân bằng năm lớp NOC. Focal Loss $\mathcal{L}_{\mathrm{focal}} = -\alpha_y(1-p_y)^{\gamma}\log p_y$ với $\gamma=2$ và trọng số lớp $\alpha_y = N_{\mathrm{total}}/(5\cdot N_y)$ tỉ lệ nghịch tần suất từng lớp.

\subsection{Lọc khả thi và tách tập peak riêng với peak chia sẻ}

Trước khi đưa vào bộ mã hóa, hai phép tiền xử lý dựa trên panel kiểu gen được áp dụng cho tập peak đã chiếu \(\bm{x}\in\mathbb{R}^{B\times N\times128}\). Từ kiểu gen tham chiếu của \(C=45\) donor, ta dựng một bảng tra mang allele \(\Omega[l,b,c]\in\{0,1\}\) cho biết donor \(c\) có mang allele rơi vào bin \(b\) của locus \(l\) hay không. Với mỗi peak, số donor mang allele của nó là
\begin{equation}
n^{\mathrm{car}}_{p}=\sum_{c=1}^{C}\Omega\!\left[l_p,\,b_p,\,c\right],
\qquad b_p=\operatorname{round}(10\,a_p)+\delta,
\end{equation}
trong đó \(a_p\) là allele của peak và \(\delta\) là độ dời chỉ số bin cố định. Bộ lọc khả thi loại khỏi tập tất cả peak mà không donor nào giải thích được, tức \(n^{\mathrm{car}}_p=0\); những peak này được đánh dấu pad và triệt về không, \(\bm{x}_p\leftarrow\bm{x}_p\,\mathbf{1}[n^{\mathrm{car}}_p>0]\). Về mặt pháp y, đây là các artefact kiểu drop-in hoặc stutter không khớp panel; loại bỏ chúng trước bộ giải mã giúp bằng chứng peak không bị quy nhầm cho một donor khi nhiệm vụ xếp hạng giả định panel đóng.

Phép thứ hai, tách tập của tập (set-of-set), chia các peak hợp lệ thành hai tập con rời nhau theo số donor mang: tập riêng gồm các peak có \(n^{\mathrm{car}}_p=1\), tức chỉ một donor duy nhất mang allele đó, và tập chia sẻ gồm các peak còn lại với \(n^{\mathrm{car}}_p\neq1\). Hai tập con đi qua \emph{cùng} một bộ mã hóa với mặt nạ riêng, rồi được cộng lại trên các giá đỡ rời nhau,
\begin{equation}
\mathbf{H}=\operatorname{Enc}(\bm{x};\,\mathcal{M}_{\text{riêng}})\odot\mathbf{1}[\text{riêng}]
+\operatorname{Enc}(\bm{x};\,\mathcal{M}_{\text{chia sẻ}})\odot\mathbf{1}[\text{chia sẻ}].
\end{equation}
Một peak đơn-donor là bằng chứng mạnh nhất để khẳng định sự hiện diện của một contributor, đặc biệt là contributor thiểu số. Khi xử lý chung với các peak chia sẻ cao của contributor chính, thống kê chuẩn hóa theo tập bị các peak cao chi phối và peak đơn-donor yếu dễ bị làm mờ; tách riêng hai tập cho phép bộ mã hóa dựng một biểu diễn sạch và đặc thù donor cho tập riêng.

Cuối cùng, vì bộ lọc khả thi loại bỏ chính bằng chứng cần cho phát hiện open-set, một nhánh từ chối song song mã hóa lại toàn bộ tập \emph{chưa lọc} trong một lượt độc lập, tách gradient rồi gộp bằng một PMA riêng để tạo điểm từ chối. Việc tách gradient bảo đảm mục tiêu open-set không làm nhiễu bộ mã hóa hay đường xếp hạng, trong khi nhánh này vẫn nhìn thấy các peak ngoài panel mà một hỗn hợp có contributor lạ để lại.

\subsection{Bộ giải mã khe có điều kiện kiểu gen}

Bộ giải mã biến nhận dạng contributor thành bài toán gán bằng chứng peak cho một tập khe, trong đó mỗi khe được neo vào một donor tham chiếu. Bốn cơ chế kế thừa từ học biểu diễn lấy vật thể làm trung tâm được điều chỉnh sang bối cảnh pháp y: khởi tạo khe có điều kiện, tinh chỉnh theo hướng dẫn, định tuyến bằng vận tải tối ưu và cổng tồn tại.

Khởi tạo có điều kiện kiểu gen theo tinh thần Conditional Slot Attention thay từ điển khe học sẵn bằng từ điển kiểu gen \cite{ma2024cosa}. Tập allele tham chiếu của donor \(c\) được mã hóa qua \emph{cùng} đường nhúng token rồi lấy trung bình, sau đó chiếu qua một lớp tuyến tính khởi tạo bằng không để tạo neo khe
\begin{equation}
\bm{g}_c=W_{\mathrm{geno}}\!\left(\frac{\sum_j m_{cj}\,\bm{e}_{cj}}{\sum_j m_{cj}}\right),
\qquad \bm{S}^{(0)}_c=\bm{g}_c\in\mathbb{R}^{128}.
\end{equation}
Khởi tạo bằng không khiến mọi khe xuất phát trung lập ở bước đầu và dịch dần về phía allele đặc trưng của donor khi huấn luyện, tránh áp đặt một thiên lệch cứng theo kiểu gen.

Bước tinh chỉnh theo hướng dẫn, lấy cảm hứng từ Guided Slot Attention \cite{lee2024gsa}, dùng một đầu attribution sinh logits gán peak cho donor \(\bm{A}^{\mathrm{attr}}\in\mathbb{R}^{B\times N\times(C+1)}\) với cột cuối là nền. Trọng số attribution chuẩn hóa được dùng để gom các peak về từng khe, sau đó một cổng học được điều tiết lượng cập nhật,
\begin{equation}
\bm{r}_c=\sum_{i=1}^{N}\operatorname{softmax}(\bm{A}^{\mathrm{attr}})_{ic}\,\mathbf{H}_i,
\qquad
\bm{S}_c\leftarrow\bm{S}_c+\sigma\!\left(w^\top W_{r}\bm{r}_c\right)W_{r}\bm{r}_c.
\end{equation}
Tín hiệu attribution đóng vai trò hướng dẫn tương thích giữa peak và kiểu gen, giúp khe của một donor được bổ sung đúng bằng chứng peak liên quan trước khi vào vòng định tuyến.

Lõi định tuyến áp dụng MESH, một dạng slot-attention bằng vận tải tối ưu \cite{zhang2023mesh}. Trong mỗi vòng, độ tương hợp giữa khe và peak là \(\bm{C}=Q(\bm{S})K(\mathbf{H})^\top/\sqrt{d}\); ma trận gán \(\bm{A}\) được tính bằng Sinkhorn trong miền log với chuẩn hóa luân phiên theo khe và theo peak,
\begin{equation}
\bm{L}\leftarrow\bm{C}/\epsilon,\qquad
\bm{L}\leftarrow\bm{L}-\operatorname*{logsumexp}_{\text{khe}}\bm{L},\qquad
\bm{L}\leftarrow\bm{L}-\operatorname*{logsumexp}_{\text{peak}}\bm{L},\qquad
\bm{A}=\exp(\bm{L}),
\end{equation}
lặp năm bước với \(\epsilon=0{,}05\). Việc chuẩn hóa hai phía buộc một peak chia sẻ được chia tỷ lệ cho nhiều khe thay vì bị một khe chiếm trọn, qua đó chống hiện tượng giải thích lấn (explaining-away) khi hai donor cùng mang nhiều allele. Khe được cập nhật qua một GRU rồi một mạng truyền thẳng,
\begin{equation}
\bm{S}\leftarrow\operatorname{GRU}\!\left(\bm{A}\,V(\mathbf{H}),\,\bm{S}\right),
\qquad
\bm{S}\leftarrow\bm{S}+\operatorname{FFN}(\operatorname{LN}(\bm{S})),
\end{equation}
với ba vòng định tuyến cho mỗi lần suy luận.

Cuối cùng, cổng tồn tại theo Adaptive Slot Attention biến số khe biến đổi thành một biến Bernoulli mềm cho mỗi donor \cite{fan2024adaslot}. Từ logit tồn tại \(e_c=w_g^\top\bm{S}_c\), cổng được lấy mẫu bằng nhiễu logistic khi huấn luyện và lấy kỳ vọng khi suy luận,
\begin{equation}
g_c=\sigma\!\left((e_c+\eta_c)/\tau\right),\quad \eta_c=\log u-\log(1-u),\ u\sim\mathcal{U}(0,1),
\end{equation}
với \(\tau=1\); nhiễu logistic là nới lỏng đúng của phân phối Bernoulli, khác với một biến Gumbel đơn. Logit contributor cộng nội dung khe và sự tồn tại trong không gian logit, tương đương nhân hai xác suất nội dung và hiện diện,
\begin{equation}
\ell_c=\operatorname{cls\_head}(\bm{S}_c)+e_c,
\qquad P(\mathrm{donor}_c\mid X)=\sigma(\ell_c).
\end{equation}

\subsection{Tái xếp hạng theo tỷ lệ hỗn hợp}

Logit contributor học được mã hóa nhiều tín hiệu trộn lẫn, nên việc bổ sung một ước lượng độc lập về tỷ lệ hỗn hợp có thể làm sắc nét thứ hạng donor. Ta ước lượng tỷ lệ này bằng một bước EM thuần dựa trên chiều cao peak và tính tương thích kiểu gen, hoàn toàn không dùng trọng số mạng, theo tinh thần mô hình liên tục kiểu EuroForMix \cite{bleka2016}. Với mỗi peak khớp một allele có donor mang, trách nhiệm của donor \(c\) và bước cập nhật tỷ lệ là
\begin{equation}
r_{ic}=\frac{\mathbf{1}[c\in\mathcal{C}_i]\,\pi_c}{\pi_{\mathrm{bg}}\,e^{-2}+\sum_{j\in\mathcal{C}_i}\pi_j},
\qquad
\pi_c\propto\sum_i r_{ic}\,h_i,
\end{equation}
trong đó \(\mathcal{C}_i\) là tập donor mang allele của peak \(i\), \(h_i\) là chiều cao quan sát và \(\pi_{\mathrm{bg}}\) là thành phần nền; chiều cao được chia cạnh tranh đều giữa các donor mang allele (tương thích đồng đều), nên kênh tỷ lệ giữ được tính độc lập với điểm của mạng. Sau mười vòng, điểm contributor được tái xếp hạng bằng một phép pool ý kiến logarit, tức cộng log-xác suất của hai chuyên gia độc lập,
\begin{equation}
R_c=z(\ell_c)+\alpha\,z\!\left(\log(\pi_c+\varepsilon)\right),
\end{equation}
với \(z(\cdot)\) là chuẩn hóa z-score theo hàng và \(\alpha\) được chọn trên tập validation, không dùng test. Cơ sở của phép cộng này là dưới giả thiết Bayes và độc lập, kết hợp hai phân phối chuyên gia tương đương với cộng log-xác suất \cite{hinton2002}; chính vì vậy tín hiệu EM tương thích đồng đều, vốn ít trùng lặp với logit, mới nâng được chất lượng xếp hạng. Trên checkpoint seed \(42\), phép tái xếp hạng nâng exact match oracle của lớp năm contributor từ \(0{,}831\) lên \(0{,}906\); vì đây là một checkpoint, kết quả được xem là bằng chứng ban đầu cần xác nhận trên nhiều seed.

\subsection{Đếm số người đóng góp và giải mã top \(k\)}

Số người đóng góp được suy ra từ chính hồ sơ xác suất donor, không cần một mạng đếm riêng. Với \(\bm{P}=\sigma(\bm{\ell})\), một bộ rừng ngẫu nhiên hậu nghiệm dự đoán \(\hat{k}\) từ một mô tả gọn của hồ sơ xác suất gồm tám xác suất lớn nhất đã sắp, tổng xác suất và số donor vượt ngưỡng \(0{,}5\). Bộ rừng được huấn luyện trên tập validation với nhãn là số người tối ưu theo chi phí
\begin{equation}
k^{*}=\arg\min_{k\in\{1,\dots,5\}}
\underbrace{\frac{\sum_c y_c\,(1-\mathrm{top}_k)_c}{\max(\sum_c y_c,1)}}_{\text{tỷ lệ bỏ sót}}
+\underbrace{\frac{\sum_c(1-y_c)\,\mathrm{top}_{k,c}}{k}}_{\text{tỷ lệ dư}}
+\lambda k,\qquad \lambda=0{,}02,
\end{equation}
trong đó \(\mathrm{top}_k\) là chỉ thị \(k\) donor xác suất cao nhất. Tập contributor cuối cùng là \(\hat{k}\) donor có điểm xếp hạng cao nhất sau bước tái xếp hạng. Khi báo cáo, ta tách \emph{oracle top true} \(k\), dùng số người thật để cô lập chất lượng xếp hạng, khỏi \emph{deployed top} \(\hat{k}\), dùng số người dự đoán để phản ánh hiệu năng triển khai.

Việc đếm được thực hiện trên hồ sơ xác suất chứ không trên điểm đã tái xếp hạng. Đếm trên điểm tái xếp hạng làm dịch chuyển lỗi giữa các lớp, đo được giảm khoảng mười ba điểm phần trăm exact match ở lớp ba và bảy điểm ở lớp bốn để đổi lấy cải thiện ở lớp năm, nên bị loại. Một đầu đếm thứ tự kiểu CORN cũng bị loại vì trên cùng tập thử nghiệm, đầu này kém hơn đáng kể đếm hậu nghiệm ở các lớp NOC cao (exact match lớp bốn và lớp năm chỉ quanh \(0{,}75\) và \(0{,}67\) so với \(0{,}93\) và \(0{,}80\) của đếm hậu nghiệm); thay vì ràng buộc thứ tự cứng như vậy, ta giữ đếm hậu nghiệm trên hồ sơ xác suất \cite{shi2023corn}. Đầu đếm năm lớp học được vẫn được giữ như một tín hiệu phụ trong huấn luyện và là hiện thực dùng cho bài toán Number of Contributors độc lập trong Hình~\ref{fig:settransv4}.

\subsection{Hàm mất mát}

Mục tiêu huấn luyện của DS-ST tập trung vào nhận dạng contributor, với các thành phần phụ định hình biểu diễn. Thành phần chính là Asymmetric Loss trên logit donor, một mở rộng của Focal Loss hạ mạnh trọng số các donor âm dễ \cite{lin2017}. Với \(p=\sigma(\ell_c)\), nhãn \(y\), biên cắt \(m\) và \(p_-=\min(1-p+m,1)\),
\begin{equation}
\mathcal{L}_{\mathrm{ASL}}=-\Big[y\,(1-p)^{\gamma_+}\log p+(1-y)\,p_-^{\,\gamma_-}\log p_-\Big],
\qquad \gamma_+=0,\ \gamma_-=4,
\end{equation}
phù hợp với việc chỉ một số ít trong \(45\) donor là dương trong mỗi hồ sơ. Điểm từ chối open-set được giám sát bằng entropy chéo nhị phân với nhãn không cho hỗn hợp khép kín và một cho hỗn hợp có contributor ngoài panel. Đầu đếm phụ học mục tiêu \(k^{*}\) bằng entropy chéo có trọng số lớp. Hai mục tiêu đặc quyền chỉ có trên dữ liệu in-silico là attribution theo peak và hồi quy tỷ lệ hỗn hợp; nhãn attribution mềm của một peak tỉ lệ với tích tỷ lệ donor và số bản sao allele rồi chuẩn hóa trên các donor mang allele,
\begin{equation}
\tilde{y}_{pc}=\frac{\pi_c\,\mathrm{CN}_{l_p,b_p,c}}{\sum_{c'}\pi_{c'}\,\mathrm{CN}_{l_p,b_p,c'}},
\end{equation}
xấp xỉ cách mô hình liên tục chia chiều cao theo kiểu gen. Hai mục tiêu phụ này được cân bằng tự động bằng trọng số bất định đồng phương sai, trong đó \(s=\log\sigma^2\) là tham số học được \cite{kendall2018}. Tổng mất mát là
\begin{equation}
\mathcal{L}=\mathcal{L}_{\mathrm{ASL}}+0{,}5\,\mathcal{L}_{\mathrm{reject}}+0{,}3\,\mathcal{L}_{\mathrm{card}}
+e^{-s_{\mathrm{attr}}}\mathcal{L}_{\mathrm{attr}}+s_{\mathrm{attr}}
+e^{-s_{\phi}}\mathcal{L}_{\phi}+s_{\phi}.
\end{equation}
Khi huấn luyện, một phép tăng cường ngẫu nhiên loại bỏ khoảng mười lăm phần trăm các peak \emph{chia sẻ} nhưng không bao giờ loại peak đơn-donor của contributor thiểu số, nhờ đó số người vẫn xác định được trong khi mô hình học được tính bền với mất peak. Đối với bài toán Number of Contributors độc lập, biến thể đếm được huấn luyện bằng Focal Loss \(\mathcal{L}_{\mathrm{focal}}=-\alpha_y(1-p_y)^{\gamma}\log p_y\) với \(\gamma=2\) cùng trọng số lớp nghịch tần suất.

\subsection{Độ đo đánh giá}

Với lớp \(k\), precision và recall được tính:
\begin{equation}
\operatorname{Precision}_k=\frac{TP_k}{TP_k+FP_k},
\qquad
\operatorname{Recall}_k=\frac{TP_k}{TP_k+FN_k}.
\end{equation}
F1 của lớp là
\begin{equation}
F1_k=\frac{2\operatorname{Precision}_k\operatorname{Recall}_k}
{\operatorname{Precision}_k+\operatorname{Recall}_k}.
\end{equation}
Macro F1 là trung bình trên năm lớp, do đó mỗi NOC có trọng số bằng nhau. Micro F1 gộp quyết định của toàn bộ mẫu và trong phân loại đơn nhãn bằng accuracy. Do phân phối NOC một chiếm hơn 80\% ở GlobalFiler, Macro F1 là chỉ số chính.

Đối với contributor identification, exact match yêu cầu toàn bộ véc tơ dự đoán bằng nhãn. Oracle top \(k\) sử dụng NOC thật để đo chất lượng ranking, còn deployed top \(k\) sử dụng NOC dự đoán. Recall at \(k\) đo tỷ lệ contributor thật nằm trong \(k\) vị trí đầu. Reject AUROC đo khả năng phân biệt mixture chỉ gồm donor trong panel với mixture có contributor ngoài panel.

\section{Thử nghiệm và đánh giá}

\subsection{Thiết lập thực nghiệm}

Các benchmark NOC được chạy trên GlobalFiler với GroupKFold năm phần; trước mỗi fold, tập huấn luyện được cân bằng bằng \texttt{balanced\_resample} — oversample lớp thiểu số và cap lớp đa số ở mức tối đa $3\times$ lớp nhỏ nhất.

\textbf{XGBoost} dùng \texttt{XGBClassifier}(n\_estimators=300, max\_depth=6, learning\_rate=0{,}05, tree\_method=\texttt{hist}, eval\_metric=\texttt{mlogloss}) trên véc tơ phẳng chiều cao allele của 24 locus GlobalFiler.

\textbf{CNN một chiều} xây dựng theo lớp tích chập \texttt{Conv1d}: khối 1 (1→64, kernel 3, BN, ReLU, \texttt{MaxPool1d(2)}), khối 2 (64→128, kernel 3, BN, ReLU, \texttt{MaxPool1d(2)}), khối 3 (128→256, kernel 3, BN, ReLU, \texttt{AdaptiveAvgPool1d(8)}). Flatten tạo véc tơ $256\times8=2048$ chiều; MLP $2048\to256\to128\to5$ với ReLU và Dropout (0{,}3 sau lớp đầu, 0{,}2 sau lớp hai). Tối ưu Adam lr=$10^{-3}$, weight decay $10^{-4}$; early stop patience 20 epoch theo Macro F1 validation.

\textbf{Forensic Transformer} tạo 24 token locus bằng cách chiếu riêng từng locus $m$ qua \texttt{Linear}($d_m\to64$) trong đó $d_m$ là số cột allele của locus đó (kích thước khác nhau giữa các locus). Chuỗi 24 token đi qua ba \texttt{TransformerEncoderLayer}(d=64, nhead=4, FFN=256, dropout=0{,}1, batch\_first=True); mean-pool trên 24 token; đầu phân loại \texttt{LayerNorm}(64)$\to$\texttt{Linear}(64$\to$5). Tối ưu Adam lr=$5\times10^{-4}$, weight decay $10^{-5}$; early stop patience 20.

\textbf{TabPFN} dùng \texttt{TabPFNClassifier}(N\_ensemble\_configurations=32); tập huấn luyện mỗi fold được lấy mẫu phân tầng tối đa 3.000 mẫu (do giới hạn context window mô hình); không có bước huấn luyện — một lần forward pass cho toàn bộ tập validation.

\textbf{DS-ST} khởi tạo từ checkpoint \texttt{noc\_10mb} (tiền huấn luyện trên bài toán contributor identification 45 lớp), dùng $d_{\mathrm{model}}=128$, bốn head, hai ISAB++, 32 inducing points, periodic embedding ($n_{\mathrm{freq}}=8$, $\sigma=0{,}3$), nc\_attn=\texttt{mab0}, feas\_filter=True, set\_of\_set=True và token tám chiều. Đầu NOC mới là Dropout(0{,}1)$\to$Linear(128$\to$64)$\to$ReLU$\to$Dropout(0{,}1)$\to$Linear(64$\to$5).

\textbf{Cấu hình tính toán.} Toàn bộ quá trình huấn luyện và đánh giá mô hình được thực hiện trên nền tảng Kaggle với một GPU NVIDIA Tesla T4 (16~GB VRAM): huấn luyện DS-ST trên 55.247 hỗn hợp bán-tổng hợp cho nhiệm vụ nhận dạng contributor, và toàn bộ benchmark Number of Contributors theo GroupKFold năm phần trên 10.124 hồ sơ GlobalFiler (XGBoost, CNN một chiều, Forensic Transformer, TabPFN và biến thể đếm DS-ST). Các mô hình học sâu được hiện thực bằng PyTorch với CUDA; XGBoost và TabPFN dùng thư viện chính thức, trong đó TabPFN nạp checkpoint tiền huấn luyện công khai và chỉ chạy in-context learning khi suy luận. Một máy tính xách tay cá nhân với GPU NVIDIA GeForce RTX 3050 Laptop (4~GB VRAM) được dùng cho các tác vụ nhẹ hơn: tiền xử lý dữ liệu, tạo hỗn hợp in-silico và phát triển, gỡ lỗi mô hình.

\subsection{Kết quả benchmark}

\begin{table}[H]
\centering
\caption{Kết quả GroupKFold năm phần trên GlobalFiler. Các mô hình dùng cùng group identity để tách fold.}
\label{tab:benchmark}
\small
\begin{tabular}{lcc}
\toprule
\textbf{Mô hình} & \textbf{Macro F1} & \textbf{Micro F1} \\
\midrule
CNN một chiều & \(0,8854\pm0,0100\) & \(0,9666\pm0,0051\) \\
Forensic Transformer & \(0,8846\pm0,0131\) & \(0,9624\pm0,0030\) \\
XGBoost & \(0,9230\pm0,0209\) & \(0,9763\pm0,0082\) \\
TabPFN & \(0,9297\pm0,0108\) & \(0,9787\pm0,0040\) \\
DS-ST & \(\mathbf{0,9444\pm0,0042}\) & \(\mathbf{0,9822\pm0,0026}\) \\
\bottomrule
\end{tabular}
\end{table}

DS-ST vượt XGBoost 0,0214 điểm Macro F1, tương đương 2,14 điểm phần trăm tuyệt đối. Khoảng cách với CNN là 0,0590 và với Forensic Transformer là 0,0598. TabPFN đạt Macro F1 0,9297, vượt XGBoost dù chỉ dùng tối đa 3.000 mẫu mỗi fold và không cần bước huấn luyện gradient, song vẫn thấp hơn DS-ST 0,0147 điểm vì không khai thác cấu trúc tập peak. Độ lệch chuẩn Macro F1 của DS-ST bằng 0,0042, nhỏ hơn XGBoost 0,0209, cho thấy kết quả ổn định hơn giữa các fold. Micro F1 cao ở mọi mô hình vì lớp NOC một chiếm đa số; chênh lệch Macro F1 phản ánh rõ hơn hiệu năng trên lớp phức tạp. Đây là kết quả của nhánh đếm DS-ST; kết quả nhận dạng contributor của cùng mô hình được trình bày ở các mục sau.

\subsection{Phân tích theo lớp của DS-ST}

\begin{table}[H]
\centering
\caption{F1 theo lớp NOC của DS-ST trên từng fold GroupKFold 5. Tất cả giá trị được trích trực tiếp từ báo cáo phân loại.}
\label{tab:perclass}
\small
\resizebox{\textwidth}{!}{%
\begin{tabular}{lccccccc}
\toprule
\textbf{Fold} & \textbf{NOC 1} & \textbf{NOC 2} & \textbf{NOC 3} & \textbf{NOC 4} & \textbf{NOC 5} & \textbf{Macro F1} & \textbf{Accuracy} \\
\midrule
1 & 0,996 & 0,939 & 0,967 & 0,934 & 0,911 & 0,9494 & 0,983 \\
2 & 0,995 & 0,967 & 0,974 & 0,917 & 0,863 & 0,9430 & 0,983 \\
3 & 0,992 & 0,894 & 0,920 & 0,934 & 0,945 & 0,9370 & 0,977 \\
4 & 0,995 & 0,917 & 0,949 & 0,958 & 0,913 & 0,9463 & 0,984 \\
5 & 0,994 & 0,943 & 0,948 & 0,895 & 0,953 & 0,9465 & 0,984 \\
\midrule
Trung bình & 0,994 & 0,932 & 0,952 & 0,928 & 0,917 & 0,9444 & 0,9822 \\
Độ lệch chuẩn & 0,001 & 0,025 & 0,019 & 0,021 & 0,032 & 0,0042 & 0,0026 \\
\bottomrule
\end{tabular}}
\end{table}

Với bảng đầy đủ năm fold, lớp NOC một ổn định nhất (F1 từ 0,992 đến 0,996, trung bình 0,994). Trong các lớp mixture, NOC ba đạt trung bình cao nhất 0,952, tiếp theo NOC hai 0,932, NOC bốn 0,928 và NOC năm 0,917. NOC năm biến động mạnh nhất, từ 0,863 ở fold hai đến 0,953 ở fold năm; NOC hai giảm xuống 0,894 ở fold ba. Dù vậy mọi lớp mixture đều đạt F1 trung bình trên 0,91, cho thấy Macro F1 cao không chỉ đến từ lớp NOC một chiếm đa số, đồng thời củng cố nhu cầu báo cáo cả Macro F1 lẫn độ lệch chuẩn giữa các fold thay vì chỉ dùng accuracy tổng thể.

\begin{table}[H]
\centering
\caption{F1 theo lớp NOC của TabPFN trên từng fold GroupKFold 5. Tập huấn luyện được cap tối đa 3.000 mẫu/fold.}
\label{tab:perclass_tabpfn}
\small
\begin{tabular}{lccccc}
\toprule
\textbf{Fold} & \textbf{NOC 1} & \textbf{NOC 2} & \textbf{NOC 3} & \textbf{NOC 4} & \textbf{NOC 5} \\
\midrule
1 & 0,995 & 0,905 & 0,949 & 0,931 & 0,876 \\
2 & 0,992 & 0,895 & 0,953 & 0,875 & 0,833 \\
3 & 0,993 & 0,921 & 0,918 & 0,907 & 0,926 \\
4 & 0,995 & 0,900 & 0,966 & 0,952 & 0,898 \\
5 & 0,996 & 0,966 & 0,890 & 0,878 & 0,931 \\
\midrule
Trung bình & 0,994 & 0,917 & 0,935 & 0,909 & 0,893 \\
\midrule
Macro F1 & \multicolumn{5}{c}{$0{,}9297\pm0{,}0108$} \\
\bottomrule
\end{tabular}
\end{table}

Đối chiếu Bảng~\ref{tab:perclass} và Bảng~\ref{tab:perclass_tabpfn} cho thấy DS-ST vượt TabPFN trên cả bốn lớp mixture theo trung bình năm fold: NOC hai \(0,932\) so với \(0,917\), NOC ba \(0,952\) so với \(0,935\), NOC bốn \(0,928\) so với \(0,909\) và NOC năm \(0,917\) so với \(0,893\). Hai mô hình ngang nhau ở lớp NOC một chiếm đa số, nên chênh lệch Macro F1 đến chủ yếu từ các lớp mixture phức tạp.

\subsection{Kết quả nhận dạng contributor}

Trên tập thử nghiệm thực gồm 3.582 hỗn hợp, ở chế độ triển khai (mô hình tự ước lượng số người), DS-ST đạt Macro F1 \(0,987\), Micro F1 \(0,986\), exact match toàn cục \(0,962\), precision \(0,977\), recall \(0,988\) và Jaccard \(0,987\); đầu từ chối open-set đạt AUROC \(0,999\) khi phân biệt 3.582 hỗn hợp khép kín với 1.366 hỗn hợp có contributor ngoài panel, còn độ chính xác đếm số người đạt \(0,967\).

Để so sánh công bằng với năm baseline cổ điển — vốn không có đầu ước lượng số người — Bảng~\ref{tab:contrib_baseline} dùng giải mã oracle top-\(k\) (\(k\) thật) cho \emph{mọi} mô hình, qua đó cô lập riêng chất lượng xếp hạng donor khỏi khâu đếm (việc ước lượng số người đã được đánh giá riêng ở benchmark NOC GroupKFold). Dưới giải mã này DS-ST vẫn dẫn đầu mọi chỉ số: exact match \(0,985\) so với \(0,956\) của XGBoost (mạnh nhất trong nhóm baseline), \(0,951\) của hồi quy logistic, \(0,949\) của MLP, \(0,924\) của CNN và \(0,780\) của kNN. Ngoài ra chỉ DS-ST có đầu từ chối open-set; các baseline đa nhãn không phân biệt được hỗn hợp có contributor ngoài panel.

\begin{table}[H]
\centering
\caption{So sánh nhận dạng contributor trên tập thử nghiệm thực (panel 45 donor, \(n=3.582\)); giải mã oracle top-\(k\) với \(k\) là số người thật cho mọi mô hình để cô lập chất lượng xếp hạng. Reject AUROC chỉ áp dụng cho mô hình có đầu từ chối.}
\label{tab:contrib_baseline}
\small
\begin{tabular}{lcccc}
\toprule
\textbf{Mô hình} & \textbf{Macro F1} & \textbf{Micro F1} & \textbf{Exact Match} & \textbf{Reject AUROC} \\
\midrule
\textbf{DS-ST} & \textbf{0,987} & \textbf{0,991} & \textbf{0,985} & \textbf{0,999} \\
XGBoost & 0,960 & 0,972 & 0,956 & --- \\
Hồi quy logistic & 0,954 & 0,968 & 0,951 & --- \\
MLP & 0,954 & 0,968 & 0,949 & --- \\
CNN & 0,936 & 0,954 & 0,924 & --- \\
kNN & 0,833 & 0,851 & 0,780 & --- \\
\bottomrule
\end{tabular}
\end{table}

\begin{table}[H]
\centering
\caption{Phân rã Exact Match theo từng mức NOC trên tập thử nghiệm thực (giải mã oracle top-\(k\), \(k\) thật, cho mọi mô hình). Cỡ mẫu mỗi lớp: 2.249 / 334 / 385 / 242 / 372.}
\label{tab:contrib_pernoc}
\small
\begin{tabular}{lcccccc}
\toprule
\textbf{Mô hình} & \textbf{NOC 1} & \textbf{NOC 2} & \textbf{NOC 3} & \textbf{NOC 4} & \textbf{NOC 5} & \textbf{Tổng} \\
\midrule
\textbf{DS-ST} & \textbf{0,998} & \textbf{0,994} & \textbf{0,995} & \textbf{0,955} & \textbf{0,906} & \textbf{0,985} \\
XGBoost & 0,996 & 1,000 & 0,961 & 0,888 & 0,710 & 0,956 \\
Hồi quy logistic & 0,993 & 0,994 & 0,964 & 0,872 & 0,696 & 0,951 \\
MLP & 0,994 & 0,994 & 0,984 & 0,855 & 0,661 & 0,949 \\
CNN & 0,988 & 0,991 & 0,930 & 0,789 & 0,554 & 0,924 \\
kNN & 0,911 & 0,787 & 0,849 & 0,364 & 0,183 & 0,780 \\
\bottomrule
\end{tabular}
\end{table}

Bảng~\ref{tab:contrib_pernoc} phân rã exact match theo từng mức NOC. Mọi mô hình gần bão hòa ở các lớp dễ (NOC một và hai), nhưng phân hóa mạnh ở lớp khó: tại NOC bốn DS-ST đạt \(0,955\) so với \(0,888\) của XGBoost và \(0,789\) của CNN, còn tại NOC năm DS-ST đạt \(0,906\) so với \(0,710\) của XGBoost và chỉ \(0,183\) của kNN. Khoảng cách giữa DS-ST và baseline vì thế nới rộng theo số người đóng góp — đúng vùng mà nhận dạng contributor trở nên khó nhất do peak của thành phần thiểu số dễ bị che lấp.

Bảng~\ref{tab:contrib} tách riêng kết quả của DS-ST theo từng mức NOC dưới hai chế độ giải mã: oracle dùng số người thật để cô lập chất lượng xếp hạng, còn post-hoc dùng số người do bộ đếm dự đoán để phản ánh hiệu năng triển khai.

\begin{table}[H]
\centering
\caption{Exact match theo NOC của DS-ST trên tập thử nghiệm thực, sau tái xếp hạng theo tỷ lệ hỗn hợp (phi-rerank). Oracle = top-true-\(k\); post-hoc = top-\(\hat{k}\) với \(\hat{k}\) do bộ đếm dự đoán trên hồ sơ xác suất.}
\label{tab:contrib}
\small
\begin{tabular}{lcccccc}
\toprule
\textbf{Chế độ} & \textbf{NOC 1} & \textbf{NOC 2} & \textbf{NOC 3} & \textbf{NOC 4} & \textbf{NOC 5} & \textbf{Tổng} \\
\midrule
\(n\) mẫu & 2.249 & 334 & 385 & 242 & 372 & 3.582 \\
EM (oracle)   & 0,998 & 0,994 & 0,995 & 0,955 & 0,906 & 0,985 \\
EM (post-hoc) & 0,997 & 0,955 & 0,917 & 0,942 & 0,815 & 0,962 \\
\bottomrule
\end{tabular}
\end{table}

Bảng~\ref{tab:contrib} đã phản ánh giải mã triển khai \emph{có} tái xếp hạng theo tỷ lệ hỗn hợp. Hiệu năng giảm đơn điệu theo số người đóng góp, đúng như kỳ vọng vì độ phức tạp tổ hợp và số contributor tỷ lệ thấp tăng theo NOC. Tái xếp hạng tác động rõ nhất ở các lớp khó: oracle NOC năm tăng từ \(0,831\) (xếp hạng thô) lên \(0,906\), còn post-hoc NOC bốn và NOC năm tăng lần lượt từ \(0,930\) và \(0,796\) lên \(0,942\) và \(0,815\); bước này chỉ đổi \emph{thứ hạng} donor chứ không động đến bộ đếm chạy trên hồ sơ xác suất. Khoảng cách còn lại giữa oracle và post-hoc nới rộng ở NOC năm (\(0,906\) so với \(0,815\)), cho thấy ở mức khó nhất lỗi đến từ \emph{cả} khâu xếp hạng \emph{và} khâu đếm.

\subsection{Quá trình phát triển từ Set Transformer đến DS-ST}

DS-ST không xuất hiện một lần mà là kết quả của một chuỗi cải tiến có kiểm soát xuất phát từ Set Transformer gốc, mỗi bước được chọn giữ hay loại dựa trên tập dev tổ hợp-mới in-silico và chỉ báo cáo một lần trên tập thử nghiệm thực. Điểm xuất phát là một Set Transformer gộp với token ba trường; trên đánh giá tổ hợp-mới, ngay cả các mô hình cổ điển có căn chỉnh như hồi quy logistic và XGBoost cũng sánh ngang hoặc vượt nó, vì phép gộp tập làm mất căn chỉnh locus\(\times\)allele mà các mô hình tra cứu theo donor tận dụng — động lực để cải tiến chính bộ mã hóa tập.

Tăng cường token lên tám trường dưới dạng scalar thô làm sụp đổ một bộ mã hóa ISAB thường (oracle giảm còn \(0,88\), NOC4 và NOC5 về gần không); thay bộ mã hóa bằng ISAB++ với SetNorm và đường residual sạch khôi phục hoàn toàn (oracle \(0,96\)), cho thấy \emph{bộ mã hóa}, chứ không phải số trường, mới là đòn bẩy chính. Thêm nhúng số tuần hoàn theo từng đặc trưng với tần số khởi tạo vừa phải (\(\sigma=0,3\)) mang lại mức tăng nhất quán trên cả dev lẫn test. Tiếp đó, giám sát đặc quyền gồm attribution theo peak và hồi quy tỷ lệ hỗn hợp, cân bằng theo bất định, nâng oracle xếp hạng và tăng gấp khoảng ba lần oracle NOC năm, đồng thời thu hẹp đáng kể khoảng cách in-silico\(\rightarrow\)thực.

Phân tích cơ chế định vị điểm yếu NOC năm còn lại ở chỗ attention softmax cạnh tranh tại bước điểm cảm ứng hút mất peak thiểu số mờ; thay nó bằng attention sigmoid không cạnh tranh là đòn bẩy đầu tiên dịch chuyển được rào cản NOC năm. Một bộ tái xếp hạng tỷ lệ hỗn hợp ký hiệu độc lập (lọc khả thi cộng EM chia mềm kiểu EuroForMix) bổ sung thêm phần oracle NOC năm và triển khai được trên mọi checkpoint. Cuối cùng, thay phép gộp per-donor bằng bộ giải mã khe có điều kiện kiểu gen (CoSA, GSANet, MESH, AdaSlot) cùng tách tập riêng/chia sẻ và bộ lọc khả thi tạo nên DS-ST hoàn chỉnh.

Nhiều đòn bẩy khác đã được thử và \emph{loại bỏ} vì không chuyển giao được sang dữ liệu thực hoặc gây mất ổn định huấn luyện: attention sparsemax và entmax (sụp đổ về gần không), bộ giải mã cộng tính và dòng kép kiểu DSMIL, các bộ chính quy tự giám sát VICReg và tái dựng, invariant-risk minimization, gộp bảo toàn khối lượng, trọng số ưu tiên thành phần thiểu số, điều kiện hóa theo bản sao và theo truy vấn kiểu gen, các biến thể khe ref-match, em-phi, noise-gate, cùng một đầu đếm thứ tự kiểu CORN. DS-ST chỉ giữ lại những thành phần có mức cải thiện chuyển giao được; bản thân kỷ luật lựa chọn này — chọn trên dev tổ hợp-mới in-silico, đánh giá trên test thực đúng một lần — là một đóng góp phương pháp luận của nghiên cứu.

\subsection{Ablation study}

Để tách đóng góp của từng thành phần, Bảng~\ref{tab:ablation} đo các cấu hình trên \emph{cùng} tập thử nghiệm thực (\(n=3.582\), một seed). Phần trên xây dựng dần từ một bộ giải mã per-donor đến DS-ST; phần dưới gắn thêm vào DS-ST những thành phần từng được cân nhắc nhưng bị loại.

\begin{table}[H]
\centering
\caption{Ablation thành phần của DS-ST trên tập thử nghiệm thực (\(n=3.582\), một seed). EM post-hoc là chỉ số triển khai; EM oracle cô lập chất lượng xếp hạng.}
\label{tab:ablation}
\small
\begin{tabular}{lccc}
\toprule
\textbf{Cấu hình} & \textbf{Macro F1} & \textbf{EM post-hoc} & \textbf{EM oracle} \\
\midrule
Bộ giải mã per-donor (nền) & 0,979 & 0,935 & 0,965 \\
\;\(+\) lọc khả thi \(+\) giám sát đặc quyền & 0,984 & 0,925 & 0,972 \\
\;\(+\) bộ giải mã khe \(+\) tách riêng/chia sẻ & 0,975 & 0,956 & 0,973 \\
\;\(+\) attention sigmoid không cạnh tranh \textbf{(DS-ST)} & \textbf{0,982} & \textbf{0,959} & \textbf{0,973} \\
\midrule
DS-ST \(+\) ref-match & 0,977 & 0,951 & 0,972 \\
DS-ST \(+\) đặc trưng em-\(\phi\) & 0,962 & 0,907 & 0,970 \\
DS-ST \(+\) noise-gate & 0,973 & 0,944 & 0,958 \\
\bottomrule
\end{tabular}
\end{table}

Hai nhóm thành phần cải thiện hai khía cạnh khác nhau. Lọc khả thi cùng giám sát đặc quyền nâng \emph{trần xếp hạng} (EM oracle \(0,965\rightarrow0,972\), Macro F1 \(0,979\rightarrow0,984\)) — tức cải thiện việc chọn \emph{ai}; ở giai đoạn per-donor, EM post-hoc thậm chí giảm nhẹ vì đầu đếm chưa theo kịp trần xếp hạng đã nâng. Bộ giải mã khe cùng phép tách tập riêng/chia sẻ sau đó nâng \emph{EM triển khai} \(0,925\rightarrow0,956\), nhờ biến các cổng tồn tại của khe thành một tín hiệu đếm hiệu chỉnh tốt hơn và cô lập peak đơn-donor của thành phần thiểu số. Attention sigmoid không cạnh tranh bổ sung phần Macro F1 và EM cuối cùng, hoàn thiện DS-ST ở mức \(0,982\) Macro F1 và \(0,959\) EM dưới giải mã thô. Để cô lập đóng góp của \emph{kiến trúc}, Bảng~\ref{tab:ablation} dùng chung một cách giải mã thô (chưa tái xếp hạng) cho mọi cấu hình; khi bật tái xếp hạng theo tỷ lệ hỗn hợp ở bước triển khai, DS-ST đạt \(0,987\) Macro F1 và \(0,962\) EM như Bảng~\ref{tab:contrib_baseline}.

Phần dưới của bảng cho thấy mỗi thành phần gắn thêm đều \emph{làm giảm} EM triển khai: ref-match \(-0,8\) điểm, noise-gate \(-1,5\) điểm và đặc trưng em-\(\phi\) \(-5,2\) điểm. Mức giảm lớn của em-\(\phi\) đến từ việc đưa tỷ lệ hỗn hợp khử chập vào làm \emph{đặc trưng đầu vào}: kênh này trở nên trùng lặp với logit và phá vỡ tính độc lập mà bộ tái xếp hạng hậu nghiệm cần, nên dùng \(\phi\) ở khâu suy luận tốt hơn nhồi vào mô hình. Kết quả này củng cố thiết kế tinh gọn của DS-ST: một thành phần chỉ được giữ khi cải thiện chỉ số triển khai, không thêm theo mặc định. Vì các giá trị trên là một seed, các chênh lệch nhỏ ở những lớp NOC hiếm nên được đọc như xu hướng, không phải kết luận điểm. Ngoài ra, tiền huấn luyện bộ mã hóa trên nhiệm vụ contributor cũng có lợi cho nhánh đếm: bỏ tiền huấn luyện làm Macro F1 của bài toán Number of Contributors (GroupKFold ba phần) giảm khoảng hai điểm phần trăm.

\subsection{So sánh với nghiên cứu trước}

\begin{table}[H]
\centering
\caption{Đối chiếu các nghiên cứu NOC. Chỉ số không được xem là so sánh trực tiếp vì dữ liệu, số lớp và protocol khác nhau.}
\label{tab:related_compare}
\small
\begin{tabularx}{\textwidth}{P{2.5cm}P{1.7cm}P{2.5cm}P{2.3cm}Y}
\toprule
\textbf{Nghiên cứu} & \textbf{Miền NOC} & \textbf{Dữ liệu} & \textbf{Kết quả báo cáo} & \textbf{Khác biệt chính} \\
\midrule
PACE \cite{marciano2017} & 1 đến 4 & 1.405 profile, 20 donor, 15 locus & khoảng 97\% test accuracy & feature engineering và Identifiler \\
RFC19 \cite{benschop2019} & 1 đến 5 & 590 PPF6C profile, 1.174 donor & khoảng 83,3\% accuracy & random forest 19 feature \\
Decision tree \cite{kruijver2021} & 1 đến 5 & 766 PROVEDIt profile & khoảng 0,77 đến 0,85 accuracy & minh bạch, phụ thuộc lọc artefact \\
TAWSEEM \cite{alamoudi2022} & 1 đến 5 & 6.000 profile, bốn multiplex & 97\% accuracy & MLP trên khoảng 6.000 feature \\
deepNoC \cite{taylor2024deepnoc} & 1 đến 10 & 100.000 profile mô phỏng & 89\% accuracy & simulator và fine tune phòng lab \\
Nghiên cứu này (DS-ST) & 1 đến 5 & GlobalFiler thực \(+\) hỗn hợp bán-tổng hợp & NOC Macro F1 0,9444; contributor EM 0,962 & một bộ mã hóa tập dùng chung; chống rò rỉ theo tổ hợp \\
\bottomrule
\end{tabularx}
\end{table}

Bảng~\ref{tab:related_compare} cho thấy accuracy cao không tự động đồng nghĩa phương pháp tốt hơn trong mọi môi trường. TAWSEEM dùng nhiều multiplex và véc tơ lớn, deepNoC dùng dữ liệu mô phỏng cho mười lớp, PACE chỉ đến bốn contributor, trong khi nghiên cứu hiện tại nhấn mạnh GroupKFold theo mixture identity và Macro F1 trên dữ liệu mất cân bằng. Một so sánh công bằng trong tương lai cần cùng dataset version, cùng threshold artefact, cùng split và cùng metric.

\subsection{Khả năng tổng quát hóa khác kit}

Identifiler Plus và PowerPlex 16 có lần lượt 3.160 và 1.022 hồ sơ. Hai tập nhỏ hơn GlobalFiler và có phân phối mất cân bằng tương tự. DS-ST xử lý khác biệt giữa các kit nhờ giữ chỉ số locus ngay trong token và che các locus không tồn tại bằng mask, thay vì cố định một lưới locus theo từng kit; do đó cùng một bộ mã hóa áp dụng được cho mọi kit mà không cần đổi chiều đầu vào. Các thử nghiệm cross-kit hoàn chỉnh (train-from-scratch, đóng băng bộ mã hóa GlobalFiler, hay fine-tune) cần được lưu dưới dạng artifact trước khi báo cáo con số chính thức; do đó bản paper này chỉ mô tả protocol và thống kê dữ liệu, không suy diễn metric chưa có bằng chứng.

\subsection{Thảo luận}

Kết quả benchmark cho thấy XGBoost vẫn rất cạnh tranh trên dữ liệu phẳng, phù hợp với nhận định rằng tree boosting thường mạnh trên dữ liệu bảng \cite{chen2016}. Tuy nhiên, DS-ST đạt Macro F1 cao hơn và ổn định hơn ở nhiệm vụ đếm, đồng thời dẫn đầu ở nhiệm vụ nhận dạng và đạt reject AUROC gần như hoàn hảo ở nhiệm vụ open-set. Ablation thành phần (Bảng~\ref{tab:ablation}) tách được đóng góp của từng khối: lọc khả thi cùng giám sát đặc quyền nâng trần xếp hạng, còn bộ giải mã khe cùng tách tập riêng/chia sẻ nâng chỉ số triển khai — hai cơ chế bổ trợ chứ không thay thế nhau.

Một ưu điểm của DS-ST đối với bối cảnh pháp y là tính khả diễn giải đến từ chính cấu trúc giải mã. Mỗi khe gắn với một donor cụ thể, nên ma trận vận tải tối ưu cho biết peak nào hỗ trợ donor nào và với khối lượng bao nhiêu; đầu attribution theo peak và ước lượng tỷ lệ hỗn hợp \(\phi\) cung cấp thêm hai lát cắt giải thích có thể đối chiếu với phán đoán chuyên gia. Đây là điều một mô hình gộp thành một véc tơ duy nhất khó cung cấp.

Việc biến donor thành slot có điều kiện kiểu gen là phép chuyển khái niệm từ học lấy vật thể làm trung tâm sang giải trộn pháp y: cả hai đều có nhiều quan sát chồng lấn cần tách về các nguồn ẩn. Tuy nhiên donor không phải vật thể ảnh và allele sharing tuân theo quy luật di truyền, nên CoSA, GSANet, MESH và AdaSlot chỉ là nguồn cảm hứng kiến trúc; kiểm định pháp y vẫn phải dựa trên hỗn hợp thực.

Một hạn chế quan trọng là các con số nhận dạng contributor và ablation thành phần hiện ở mức một seed, trong khi benchmark Number of Contributors dùng GroupKFold năm phần. Các chênh lệch nhỏ ở những lớp NOC hiếm vì thế nên được đọc như xu hướng. Phiên bản tiếp theo cần lặp nhiều seed cho nhánh contributor và một benchmark cross-kit có artifact đầy đủ.

\subsection{Các đóng góp chính của bài báo}

Kết quả thực nghiệm hỗ trợ các đóng góp chính. Một bộ mã hóa tập duy nhất phục vụ hai nhiệm vụ chính: ở nhận dạng (CLS), DS-ST đạt exact match 0,962 và Macro F1 0,987 trên tập thử nghiệm thực, dẫn đầu các baseline cổ điển; ở từ chối open-set (reject), đầu từ chối đạt AUROC 0,999. Bài toán đếm (NOC), tuy chỉ là đại lượng suy ra từ hồ sơ xác suất CLS, lại đạt Macro F1 0,9444 (cao hơn XGBoost 0,9230 dưới cùng GroupKFold năm phần) — một kết quả phụ đáng chú ý phát hiện trong quá trình phát triển. Ablation thành phần cho thấy giám sát đặc quyền nâng trần xếp hạng còn bộ giải mã khe nâng chỉ số triển khai, và các thành phần bị loại (ref-match, em-\(\phi\), noise-gate) đều làm giảm hiệu năng. Cuối cùng, quy trình đánh giá chống rò rỉ theo tổ hợp người đóng góp — train bán-tổng hợp, test thực với tổ hợp tách biệt — bảo đảm các con số phản ánh khả năng tổng quát hóa thật.

\section{Kết luận và hướng phát triển}

\subsection{Kết luận}

Nghiên cứu đã xây dựng một mô hình học sâu, Donor-Slot Set Transformer, cho hai nhiệm vụ trung tâm của phân tích mixture DNA STR: phân loại người đóng góp trong panel (CLS) và từ chối open-set (reject). Điểm cốt lõi là xem hồ sơ như một tập peak bất biến hoán vị và xử lý bằng một bộ mã hóa tập duy nhất. Token tám trường giàu thông tin cùng nhúng số tuần hoàn mã hóa từng peak, hai khối ISAB++ với attention sigmoid không cạnh tranh học quan hệ giữa mọi peak mà không làm mờ peak yếu của contributor thiểu số; từ biểu diễn dùng chung, bộ giải mã khe có điều kiện kiểu gen với định tuyến vận tải tối ưu phân bổ bằng chứng peak về từng donor cho CLS, và một đầu từ chối tách rời cho reject. Số người đóng góp được suy ra từ hồ sơ xác suất CLS bằng một bộ đếm hậu nghiệm; đại lượng phái sinh này được phát hiện đạt độ chính xác cao trong quá trình phát triển và được báo cáo như một kết quả phụ.

Ở bài toán Number of Contributors trên 10.124 hồ sơ GlobalFiler thuộc 3.418 mixture group, DS-ST đạt Macro F1 \(0,9444\pm0,0042\), cao hơn XGBoost \(0,9230\pm0,0209\), TabPFN \(0,9297\pm0,0108\), CNN \(0,8854\pm0,0100\) và Forensic Transformer \(0,8846\pm0,0131\). Ở bài toán nhận dạng contributor trên 3.582 hỗn hợp thực, DS-ST đạt exact match \(0,962\), Macro F1 \(0,987\) và reject AUROC \(0,999\), vượt các baseline XGBoost, MLP, hồi quy logistic, CNN và kNN. Ablation thành phần xác nhận đóng góp riêng của giám sát đặc quyền, bộ giải mã khe và attention sigmoid không cạnh tranh, đồng thời cho thấy các thành phần bị loại đều làm giảm hiệu năng.

Đối với contributor identification, bộ giải mã genotype conditioned slot cung cấp một cách diễn giải trực tiếp: mỗi slot ứng với một donor và ma trận optimal transport thể hiện peak nào hỗ trợ donor nào. Phi reranking seed 42 làm oracle exact match NOC năm tăng từ 0,831 lên 0,906, nhưng cần đánh giá nhiều seed trước khi coi mức tăng là ổn định.

\subsection{Hướng phát triển}

Hướng phát triển gần nhất là lặp lại ablation thành phần của DS-ST trên nhiều seed và cùng một tập tách, để biến các xu hướng một-seed hiện tại thành khoảng tin cậy; ưu tiên các khối còn ở mức bằng chứng một-seed gồm attention sigmoid không cạnh tranh, bộ giải mã khe và bước tái xếp hạng theo tỷ lệ hỗn hợp.

Một hướng riêng cho bài toán đếm: tuy đại lượng suy ra này đã đạt độ chính xác cao, nó vẫn còn dư địa tối ưu. Hiện tại bước tái xếp hạng theo tỷ lệ hỗn hợp (logarithmic opinion pool) chỉ được dùng để xếp hạng donor, còn bộ đếm triển khai chạy trên hồ sơ xác suất bằng rừng ngẫu nhiên — bởi vì nếu đếm trực tiếp trên điểm đã tái xếp hạng thì cải thiện được lớp NOC năm nhưng lại đánh đổi mất lớp NOC ba và NOC bốn. Vì thế, việc tối ưu để một cách tái xếp hạng thống nhất nâng được NOC năm mà không hy sinh NOC ba/bốn, qua đó đếm và xếp hạng dùng chung một điểm, là một hướng mở đáng theo đuổi.

Hướng thứ hai là huấn luyện đa nhiệm end to end. NOC head và contributor head hiện chia sẻ backbone về mặt kiến trúc nhưng có thể được tối ưu theo giai đoạn. Một loss consistency có thể ép tổng gate contributor gần kỳ vọng NOC:
\begin{equation}
\mathcal{L}_{\mathrm{cons}}=
\left|
\sum_{c=1}^{C}\sigma(e_c)
-\sum_{k=1}^{5}k\,p_k^{\mathrm{noc}}
\right|.
\end{equation}
Thiết kế này cho phép hai nhánh hỗ trợ nhau nhưng cần kiểm soát để lỗi đếm không làm hỏng ranking donor.

Hướng thứ ba là đánh giá cross kit có kiểm soát trên Identifiler Plus và PowerPlex 16. Cần so sánh train from scratch, frozen GlobalFiler encoder, fine tune toàn bộ và adapter theo kit. Locus không tồn tại phải giữ mask, còn locus chung phải giữ cùng chỉ số. Báo cáo cần có Macro F1, calibration, per class F1 và performance theo mixture ratio.

Hướng thứ tư là tích hợp allele calling từ electropherogram thô. DNANet cho thấy U Net có thể đạt F1 0,962 trên allele donor thật \cite{dewit2025}. Một pipeline kết hợp DNANet với DS-ST có thể truyền uncertainty ở mức scan point tới peak và toàn hồ sơ, thay vì coi peak extraction là hoàn toàn chính xác.

Hướng cuối cùng là phát triển explainability phù hợp pháp y. Attention map không tự động là explanation. Cần kết hợp counterfactual profile hợp lệ như Veldhuis và cộng sự, kiểm tra sensitivity theo peak dropout, stutter và thay đổi mixture proportion, đồng thời đánh giá explanation với chuyên gia \cite{veldhuis2022}. Mục tiêu cuối không chỉ là tăng metric mà là tạo hệ thống có thể kiểm chứng, tái lập và sử dụng như công cụ hỗ trợ thay vì thay thế phán đoán chuyên môn.

\newpage
\begin{thebibliography}{99}

\bibitem{jeffreys1985}
Jeffreys, A. J., Wilson, V., \& Thein, S. L. (1985).
Individual-specific fingerprints of human DNA.
\textit{Nature}, 316, 76--79.
\url{https://doi.org/10.1038/316076a0}

\bibitem{nist2024}
National Institute of Standards and Technology. (2024).
\textit{DNA Mixture Interpretation: A Scientific Foundation Review}.
NIST IR 8351.
\url{https://doi.org/10.6028/NIST.IR.8351}

\bibitem{provedit2018}
Alfonse, L. E., Garrett, A. D., Lun, D. S., Duffy, K. R., \& Grgicak, C. M. (2018).
A large-scale dataset of single and mixed-source short tandem repeat profiles to inform human identification strategies: PROVEDIt.
\textit{Forensic Science International: Genetics}, 32, 62--70.
\url{https://doi.org/10.1016/j.fsigen.2017.10.006}

\bibitem{proveditweb}
Laboratory for Forensic Technology Development and Integration. (2026).
PROVEDIt Database.
\url{https://lftdi.camden.rutgers.edu/provedit/}

\bibitem{haned2011}
Haned, H., Pène, L., Lobry, J. R., Dufour, A. B., \& Pontier, D. (2011).
Estimating the number of contributors to forensic DNA mixtures: Does maximum likelihood perform better than maximum allele count?
\textit{Journal of Forensic Sciences}, 56(1), 23--28.
\url{https://doi.org/10.1111/j.1556-4029.2010.01550.x}

\bibitem{alfonse2017}
Alfonse, L. E., Tejada, G., Swaminathan, H., Lun, D. S., \& Grgicak, C. M. (2017).
Inferring the number of contributors to complex DNA mixtures using three methods: Exploring the limits of low-template DNA interpretation.
\textit{Journal of Forensic Sciences}, 62(2), 308--316.
\url{https://doi.org/10.1111/1556-4029.13284}

\bibitem{grgicak2020}
Grgicak, C. M., Karkar, S., Yearwood-Garcia, X., Alfonse, L. E., Duffy, K. R., \& Lun, D. S. (2020).
A large-scale validation of NOCIt's a posteriori probability of the number of contributors and its integration into forensic interpretation pipelines.
\textit{Forensic Science International: Genetics}, 47, 102296.
\url{https://doi.org/10.1016/j.fsigen.2020.102296}

\bibitem{marciano2017}
Marciano, M. A., \& Adelman, J. D. (2017).
PACE: Probabilistic Assessment for Contributor Estimation, a machine learning-based assessment of the number of contributors in DNA mixtures.
\textit{Forensic Science International: Genetics}, 27, 82--91.
\url{https://doi.org/10.1016/j.fsigen.2016.11.006}

\bibitem{benschop2019}
Benschop, C. C. G., van der Linden, J., Hoogenboom, J., Ypma, R. J. F., \& Haned, H. (2019).
Automated estimation of the number of contributors in autosomal short tandem repeat profiles using a machine learning approach.
\textit{Forensic Science International: Genetics}, 43, 102150.
\url{https://doi.org/10.1016/j.fsigen.2019.102150}

\bibitem{kruijver2021}
Kruijver, M., Kelly, H., Cheng, K., Lin, M. H., Morawitz, J., Russell, L., Buckleton, J., \& Bright, J. A. (2021).
Estimating the number of contributors to a DNA profile using decision trees.
\textit{Forensic Science International: Genetics}, 50, 102407.
\url{https://doi.org/10.1016/j.fsigen.2020.102407}

\bibitem{alamoudi2022}
Alamoudi, E., Mehmood, R., Albeshri, A., \& Gojobori, T. (2022).
TAWSEEM: A deep-learning-based tool for estimating the number of unknown contributors in DNA profiling.
\textit{Electronics}, 11(4), 548.
\url{https://doi.org/10.3390/electronics11040548}

\bibitem{veldhuis2022}
Veldhuis, M. S., Ariëns, S., Ypma, R. J. F., Abeel, T. E. P. M. F., \& Benschop, C. C. G. (2022).
Explainable artificial intelligence in forensics: Realistic explanations for number of contributor predictions of DNA profiles.
\textit{Forensic Science International: Genetics}, 56, 102632.
\url{https://doi.org/10.1016/j.fsigen.2021.102632}

\bibitem{taylor2024deepnoc}
Taylor, D., \& Humphries, M. A. (2024).
deepNoC: A deep learning system to assign the number of contributors to a short tandem repeat DNA profile.
\textit{arXiv preprint arXiv:2412.09803}.
\url{https://doi.org/10.48550/arXiv.2412.09803}

\bibitem{barash2024}
Barash, M., McNevin, D., Fedorenko, V., \& Giverts, P. (2024).
Machine learning applications in forensic DNA profiling: A critical review.
\textit{Forensic Science International: Genetics}, 69, 102994.
\url{https://doi.org/10.1016/j.fsigen.2023.102994}

\bibitem{phan2021}
Phan, N. N., Chattopadhyay, A., Lee, T. T., Yin, H. I., Lu, T. P., Lai, L. C., Hwa, H. L., Tsai, M. H., \& Chuang, E. Y. (2021).
High-performance deep learning pipeline predicts individuals in mixtures of DNA using sequencing data.
\textit{Briefings in Bioinformatics}, 22(6), bbab283.
\url{https://doi.org/10.1093/bib/bbab283}

\bibitem{dewit2025}
de Wit, A. K. J. G., Wagenaar, C. D., Janssen, N. A. C., Hoegen, B., van de Wetering, J., Hoofs, H., Ariëns, S., Benschop, C. C. G., \& Ypma, R. J. F. (2026).
Making AI accessible for forensic DNA profile analysis.
\textit{Forensic Science International: Genetics}, 81, 103345. Published online August 19, 2025.
\url{https://doi.org/10.1016/j.fsigen.2025.103345}

\bibitem{bleka2016}
Bleka, Ø., Storvik, G., \& Gill, P. (2016).
EuroForMix: An open source software based on a continuous model to evaluate STR DNA profiles from a mixture of contributors with artefacts.
\textit{Forensic Science International: Genetics}, 21, 35--44.
\url{https://doi.org/10.1016/j.fsigen.2015.11.008}

\bibitem{bright2016}
Bright, J. A., Taylor, D., McGovern, C., Cooper, S., Russell, L., Abarno, D., \& Buckleton, J. (2016).
Developmental validation of STRmix, expert software for the interpretation of forensic DNA profiles.
\textit{Forensic Science International: Genetics}, 23, 226--239.
\url{https://doi.org/10.1016/j.fsigen.2016.05.007}

\bibitem{coble2019}
Coble, M. D., \& Bright, J. A. (2019).
Probabilistic genotyping software: An overview.
\textit{Forensic Science International: Genetics}, 38, 219--224.
\url{https://doi.org/10.1016/j.fsigen.2018.11.009}

\bibitem{vaswani2017}
Vaswani, A., Shazeer, N., Parmar, N., Uszkoreit, J., Jones, L., Gomez, A. N., Kaiser, Ł., \& Polosukhin, I. (2017).
Attention is all you need.
\textit{Advances in Neural Information Processing Systems}, 30, 5998--6008.
\url{https://proceedings.neurips.cc/paper/7181-attention-is-all-you-need}

\bibitem{zaheer2017}
Zaheer, M., Kottur, S., Ravanbakhsh, S., Póczos, B., Salakhutdinov, R., \& Smola, A. J. (2017).
Deep Sets.
\textit{Advances in Neural Information Processing Systems}, 30, 3391--3401.
\url{https://papers.nips.cc/paper/6931-deep-sets}

\bibitem{lee2019}
Lee, J., Lee, Y., Kim, J., Kosiorek, A., Choi, S., \& Teh, Y. W. (2019).
Set Transformer: A framework for attention-based permutation-invariant neural networks.
\textit{Proceedings of the 36th International Conference on Machine Learning}, 97, 3744--3753.
\url{https://proceedings.mlr.press/v97/lee19d.html}

\bibitem{zhang2023setnorm}
Zhang, L., Tozzo, V., Higgins, J. M., \& Ranganath, R. (2022).
Set norm and equivariant skip connections: Putting the deep in deep sets.
\textit{Proceedings of the 39th International Conference on Machine Learning}, 162, 26559--26574.
\url{https://proceedings.mlr.press/v162/zhang22ac.html}

\bibitem{gorishniy2022}
Gorishniy, Y., Rubachev, I., \& Babenko, A. (2022).
On embeddings for numerical features in tabular deep learning.
\textit{Advances in Neural Information Processing Systems}, 35, 24991--25004.
\url{https://proceedings.neurips.cc/paper_files/paper/2022/hash/9e9f0ffc3d836836ca96cbf8fe14b105-Abstract-Conference.html}

\bibitem{hendrycks2016}
Hendrycks, D., \& Gimpel, K. (2016).
Gaussian error linear units.
\textit{arXiv preprint arXiv:1606.08415}.
\url{https://doi.org/10.48550/arXiv.1606.08415}

\bibitem{lin2017}
Lin, T. Y., Goyal, P., Girshick, R., He, K., \& Dollár, P. (2017).
Focal Loss for Dense Object Detection.
\textit{Proceedings of the IEEE International Conference on Computer Vision}, 2980--2988.
\url{https://doi.org/10.1109/ICCV.2017.324}

\bibitem{shi2023corn}
Shi, X., Cao, W., \& Raschka, S. (2023).
Deep neural networks for rank-consistent ordinal regression based on conditional probabilities.
\textit{Pattern Analysis and Applications}, 26, 941--955.
\url{https://doi.org/10.1007/s10044-023-01181-9}

\bibitem{chen2016}
Chen, T., \& Guestrin, C. (2016).
XGBoost: A scalable tree boosting system.
\textit{Proceedings of the 22nd ACM SIGKDD International Conference on Knowledge Discovery and Data Mining}, 785--794.
\url{https://doi.org/10.1145/2939672.2939785}

\bibitem{hollmann2023}
Hollmann, N., Müller, S., Eggensperger, K., \& Hutter, F. (2023).
TabPFN: A Transformer that solves small tabular classification problems in a second.
\textit{International Conference on Learning Representations}.
\url{https://openreview.net/forum?id=eu9fVjVasr4}

\bibitem{locatello2020slot}
Locatello, F., Weissenborn, D., Unterthiner, T., Mahendran, A., Heigold, G., Uszkoreit, J., Dosovitskiy, A., \& Kipf, T. (2020).
Object-centric learning with Slot Attention.
\textit{Advances in Neural Information Processing Systems}, 33, 11525--11538.
\url{https://proceedings.neurips.cc/paper/2020/hash/8511df98c02ab60aea1b2356c013bc0f-Abstract.html}

\bibitem{ma2024cosa}
Kori, A., Locatello, F., De Sousa Ribeiro, F., Toni, F., \& Glocker, B. (2024).
Grounded object-centric learning.
\textit{International Conference on Learning Representations}.
\url{https://openreview.net/forum?id=pBxeZ6pVUD}

\bibitem{lee2024gsa}
Lee, M., Cho, S., Lee, D., Park, C., Lee, J., \& Lee, S. (2024).
Guided Slot Attention for unsupervised video object segmentation.
\textit{Proceedings of the IEEE/CVF Conference on Computer Vision and Pattern Recognition}.
\url{https://doi.org/10.48550/arXiv.2303.08314}

\bibitem{zhang2023mesh}
Zhang, Y., Zhang, D. W., Lacoste-Julien, S., Burghouts, G. J., \& Snoek, C. G. M. (2023).
Unlocking Slot Attention by changing optimal transport costs.
\textit{Proceedings of the 40th International Conference on Machine Learning}, 202, 41931--41951.
\url{https://proceedings.mlr.press/v202/zhang23ba.html}

\bibitem{fan2024adaslot}
Fan, K., Bai, Z., Xiao, T., He, T., Horn, M., Fu, Y., Locatello, F., \& Zhang, Z. (2024).
Adaptive Slot Attention: Object discovery with dynamic slot number.
\textit{Proceedings of the IEEE/CVF Conference on Computer Vision and Pattern Recognition}, 23062--23071.
\url{https://openaccess.thecvf.com/content/CVPR2024/html/Fan_Adaptive_Slot_Attention_Object_Discovery_with_Dynamic_Slot_Number_CVPR_2024_paper.html}

\bibitem{hinton2002}
Hinton, G. E. (2002).
Training products of experts by minimizing contrastive divergence.
\textit{Neural Computation}, 14(8), 1771--1800.
\url{https://doi.org/10.1162/089976602760128018}

\bibitem{kendall2018}
Kendall, A., Gal, Y., \& Cipolla, R. (2018).
Multi-task learning using uncertainty to weigh losses for scene geometry and semantics.
\textit{Proceedings of the IEEE Conference on Computer Vision and Pattern Recognition}, 7482--7491.
\url{https://doi.org/10.1109/CVPR.2018.00781}

\end{thebibliography}

\end{document}
